
\documentclass[11pt]{article}

\usepackage[margin=1in]{geometry}
\usepackage{amsmath,amssymb,amsthm}
\usepackage{array}
\usepackage[T1]{fontenc}
\usepackage{lmodern}
\usepackage[hidelinks]{hyperref}

\newtheorem{theorem}{Theorem}
\newtheorem{lemma}{Lemma}
\newtheorem{proposition}{Proposition}
\newtheorem{corollary}{Corollary}
\newtheorem{heuristic}{Heuristic}
\newtheorem{remark}{Remark}
\newtheorem{definition}{Definition}
\newtheorem{conjecture}{Conjecture}

\newcommand{\F}{\mathbb F}
\newcommand{\RS}{\operatorname{RS}}
\newcommand{\Bnum}{B_{\rm num}}
\newcommand{\Deltaone}{\Delta_1}
\newcommand{\Deltazero}{\Delta_0}
\newcommand{\tauone}{\tau_1}
\newcommand{\tautwo}{\tau_2}
\newcommand{\tauthree}{\tau_3}
\newcommand{\Ctwo}{C_2}
\newcommand{\qline}{q_{\rm line}}
\newcommand{\qgen}{q_{\rm gen}}
\newcommand{\qchal}{q_{\rm chal}}
\newcommand{\MCA}{\operatorname{MCA}}
\newcommand{\CA}{\operatorname{CA}}
\newcommand{\List}{\operatorname{List}}
\newcommand{\PG}{\operatorname{PG}}
\newcommand{\BCHKS}{\operatorname{BCHKS}}
\newcommand{\dist}{\Delta}
\newcommand{\efib}{c_{\rm fib}}
\newcommand{\ellpg}{\ell_{\rm pg}}
\newcommand{\epsstar}{\varepsilon_*}
\newcommand{\epsmca}{\varepsilon_{\rm mca}}
\newcommand{\epsca}{\varepsilon_{\rm ca}}
\newcommand{\Fib}{\operatorname{Fib}}
\newcommand{\Lam}{\Lambda}
\newcommand{\Bad}{\operatorname{Bad}}
\newcommand{\LD}{\operatorname{LD}}
\newcommand{\Htwo}{H_2}
\newcommand{\QH}{\mathcal Q_H}

\title{Experiments for Smooth RS-MCA Proximity Thresholds\\
\large Fresh-Eyes Threshold Memo, Missing Theorems, and Retained Locator Ledger}
\author{RS-MCA experimental notes}
\date{June 26, 2026}

\begin{document}
\maketitle

\begin{abstract}
This replacement starts with the current ``fresh eyes'' synthesis for the
smooth Reed--Solomon Proximity Prize problem.  The negative side is now quite
rigid: Paper D v5 gives a prize-facing upper cap on the MCA threshold through a
self-contained deep-point route, while Paper A and the list lower bounds rule
out any no-slack capacity theorem.  The positive side is not yet a theorem; it
should be organized around two local-limit problems, L1 arbitrary-word locator
fibers and M1 all-line residue packing, plus extension-line MCA, Paper D
conversion/radius audits, and exact protocol-ledger rewrites.

The main new additions over the old \texttt{experiments.tex} are two compact
bridge lemmas.  First, for ordinary column-distance interleaving, interleaved
lists inject into common feasible supports, so a uniform base locator-fiber
theorem gives the same polynomial exponent for all protocol arities; there is no
automatic \(n^{\mu B}\) product loss.  Second, a list bound plus an MCA bound
implies a concrete line-decoding bound with loss \(M+nL+1\).  The older BCHKS
provenance notes, quotient-locator dictionaries, normal forms, arithmetic gates,
and finite-threshold infrastructure are retained below whenever they still look
useful.
\end{abstract}

\tableofcontents

\section{Fresh-Eyes Settlement Memo}
\label{sec:fresh-eyes-settlement-memo}

I would frame the current state as follows.  The attached suite substantially
settles the negative side and gives a plausible corrected theorem shape, but it
does not yet settle the Proximity Prize threshold in the sharp ``largest
\(\delta^*\)'' sense.  The clean prize-relevant statement is no longer
``prove smooth RS reaches capacity''.  It is
\[
  \delta=1-\rho-\eta,
  \qquad
  \eta\text{ must clear generated-field entropy, quotient-core, MCA-floor,}
\]
\[
  \text{interleaving, line-field, and challenge-field budgets.}
\]
For
\[
  C=\RS[\F,D,k],\qquad n=|D|,
  \qquad \rho=k/n\in\{1/2,1/4,1/8,1/16\},
\]
the official challenge statements use this MCA/list-threshold form; in this
experimental ledger the concrete prize envelope is~\cite{ProximityPrize26}
\[
  \epsstar=2^{-128},\qquad k\le2^{40},\qquad |\F|<2^{256}.
\]
The protocol-relevant soundness contribution has the schematic form
\[
  \epsmca(C,\delta)
  +\frac{|\Lam(C^{\equiv\mu},\delta)|}{\qline}
  +\text{query error}.
\]
Thus a list theorem alone is not enough unless it is connected to the exact MCA,
correlated-agreement, line-decoding, or curve-MCA quantity used by the protocol,
and unless the denominator is the actual line or challenge field.

\subsection{Bottom line}

For the \textbf{MCA/correlated-agreement threshold}, the strongest attached
result is Paper D v5's universal cap.  Through the self-contained deep-point
route and the stated dyadic-divisor hypotheses,
\[
  \delta^*_{\rm MCA}(2^{-128})
  \le 1-\rho-2^{-9}
  \quad (\rho=1/2,1/4,1/8),
\]
and
\[
  \delta^*_{\rm MCA}(2^{-128})
  \le 1-\rho-2^{-10}
  \quad (\rho=1/16).
\]
Equivalently, throughout the relevant unsafe intervals Paper D gives
\[
  \epsmca(C,\delta)
  >\frac1{2k}\left(1-\frac n{|\F|}\right)
  \ge2^{-86}\gg2^{-128}.
\]
This is a real prize-facing cap, not merely a heuristic obstruction.  It is not
an error-one theorem; it proves failure probability on the order of \(1/k\),
which is already far above \(2^{-128}\), but it does not characterize the full
failure curve.

For the \textbf{interleaved-list threshold}, Paper A's base list lower bound
plus the base-to-interleaved monotonicity gives a clean negative cap.  At quotient
scale \(N\mid n\), with \(\rho N\in\mathbb Z\), there is a received word with at
least
\[
  \binom{N}{\rho N+1}/q,
  \qquad q=|\F|,
\]
base-code codewords at radius \(1-\rho-1/N\).  Therefore the same radius is
unsafe for the interleaved-list budget whenever this exceeds \(2^{-128}q\).  In
the worst case \(q<2^{256}\), the following choices suffice:
\[
\begin{array}{c|c|c|c|c}
\hline
\rho & N & \text{list-failure radius} &
\log_2\binom{N}{\rho N+1}-256 &
\text{margin over }2^{128}\\
\hline
1/2  & 512  & 1/2-2^{-9}    & 251.17 & 123.17\\
1/4  & 512  & 3/4-2^{-9}    & 156.33 & 28.33\\
1/8  & 1024 & 7/8-2^{-10}   & 298.68 & 170.68\\
1/16 & 2048 & 15/16-2^{-11} & 433.89 & 305.89\\
\hline
\end{array}
\]
The fourth column is the exponent of the lower-bound list size after dividing by
the worst-case field size \(2^{256}\); the fifth subtracts the target exponent
\(128\).

\subsection{What is already proved or very close to proved}

The old no-slack smooth-domain capacity claim is dead.  The quotient-locator
construction gives explicit bad slopes for canonical lines and, over deployed
smooth prime fields such as BabyBear, KoalaBear, and \(3\cdot2^{30}+1\), yields
support-wise line-MCA error one on intervals starting at
\(1-\rho-2^{-18}\) for all four prize rates, with the sharper \(2^{-17}\) gap
at rates \(1/2\) and \(1/4\).  These statements rule out any theorem saying
that every radius below \(1-\rho\) has negligible MCA error.

Paper D is the most prize-relevant negative theorem because it converts a deep
point/list phenomenon into MCA at the exact \(2^{-128}\)-style threshold.  In
v5, the headline MCA cap is routed locally through a self-contained deep-point
argument rather than the Crites--Stewart list-to-agreement import.  The remaining
CA/list comparison and fallback statements still need their own radius,
normalization, augmented-code, and import audits.

The list lower-bound side is also stronger than it is currently advertised.  It
should be promoted to an explicit interleaved-list negative corollary: base lists
embed into ordinary interleaved lists, and the binomial numerators above already
exceed the \(2^{-128}|\F|\) budget in the full field envelope for the displayed
quotient divisors.

\subsection{The corrected positive theorem shape}

The positive theorem should use
\[
  a=\lceil(1-\delta)n\rceil=k+\sigma,
  \qquad \delta=1-\rho-\eta,
\]
and require at least
\[
  \sigma\log_2\qgen
  \ge (1+\varepsilon)\log_2\binom n{k+\sigma},
  \qquad
  \sigma\ge C\frac n{\log n},
\]
plus explicit quotient-profile and failure-floor ledgers.  The first inequality
uses the generated field \(\qgen\), not a larger extension or challenge field
that merely appears later in the protocol.  The MCA numerator must be divided by
the actual line field \(\qline\).

The intended list target is an arbitrary-word locator-fiber theorem
\[
  \sup_U |\Fib_U(k+\sigma)|\le n^B.
\]
The intended MCA target is an all-line bound of the form
\[
  \epsmca(C,1-\rho-\eta)
  \le
  \frac{
    n^{1+o(1)}
    +2^{(\beta(\rho)/\Htwo(\rho))\QH(\eta)(1+o(1))}
  }{\qline},
\]
where quotient-periodic classes are removed or charged by \(\QH(\eta)\).  The
current packet contains important special cases and reductions, but not these
full arbitrary-word and all-line theorems.

\subsection{What remains missing}

\begin{enumerate}
\item \textbf{L1: arbitrary-word locator local limit.}  This is the real list
upper-bound theorem.  Characteristic-zero rigidity, monomial-prefix lanes, and
split-prime transfer results are useful but do not yet imply the uniform
polynomial-field theorem needed for arbitrary received words.

\item \textbf{M1: all-line aperiodic residue packing.}  Canonical-line exactness
and quotient floors are lower-bound tools, not all-line upper bounds.  The
positive MCA theorem needs a fixed-field, all-line, all-base residue-packing
bound.

\item \textbf{F1: extension-line MCA.}  The list side has a clean extension-code
identity.  MCA does not.  Base-valued bad lines can deflate when sampled over an
extension field, while Paper D's cap indicates that genuinely extension-valued
bad lines exist in relevant regimes.  Any extension-line lift must therefore be
a corrected-reserve theorem, not a blanket subfield statement.

\item \textbf{Paper D conversion audit.}  Paper D v5's headline MCA cap is no
longer load-bearing on the direct Crites--Stewart conversion.  The self-contained
deep-point step, support-wise monotonicity, integer-radius conventions, and the
remaining CA/list import statements should still be independently checked.

\item \textbf{Protocol-ledger rewrite.}  Deployed FRI/WHIR/IOP analyses must be
rewritten with exact terms
\[
  |\Lam(C^{\equiv\mu},\delta)|/\qline
  \quad\text{and}\quad
  \epsmca\text{ or line-decoding over the actual line field}.
\]
The survey toy protocol already has this shape; deployed analyses still need the
same ledger.
\end{enumerate}

\section{New Bridge I: Common-Support Interleaved Fiber Injection}
\label{sec:common-support-interleaved-fiber-injection}

\textbf{Status: \textsc{New-local} wrapper.}
The older ledger treated interleaving conservatively through a safe product
bound \(L_\mu(\delta)\le L_1(\delta)^\mu\).  For ordinary column-distance RS
interleaving, the following support-level injection gives the sharper upper
bound: once the one-row locator-fiber theorem is proved, interleaving costs no
extra exponent.

For \(S\subseteq D\), let
\[
  L_S(X)=\prod_{x\in S}(X-x).
\]
For a received word \(U\in\F^D\), the condition \(\deg(U\bmod L_S)<k\) means
that the unique polynomial of degree less than \(|S|\) interpolating \(U|_S\)
has degree less than \(k\).

\begin{definition}[one-row and common interleaved fibers]
Let \(C=\RS[\F,D,k]\), let \(a\ge k\), and let
\(\mathbf U=(U_1,\ldots,U_\mu)\in(\F^\mu)^D\).  Define
\[
  \Fib^{(\mu)}_{\mathbf U}(a)
  =
  \{S\subseteq D:
    |S|=a,\ \deg(U_i\bmod L_S)<k
    \text{ for every }i=1,\ldots,\mu\}.
\]
For one row, write \(\Fib_{U_i}(a)\) for the same condition with only \(U_i\).
\end{definition}

\begin{lemma}[common-support interleaved fiber injection]
\label{lem:common-support-interleaved-fiber-injection}
Let \(C=\RS[\F,D,k]\), \(n=|D|\), and let
\(a=\lceil(1-\delta)n\rceil\ge k\).  For every received \(\mu\)-row word
\(\mathbf U\in(\F^\mu)^D\), ordinary column-distance interleaving satisfies
\[
  |\Lam(C^{\equiv\mu},\delta,\mathbf U)|
  \le |\Fib^{(\mu)}_{\mathbf U}(a)|
  \le \min_{1\le i\le\mu}|\Fib_{U_i}(a)|.
\]
Consequently,
\[
  \sup_{\mathbf U}|\Lam(C^{\equiv\mu},\delta,\mathbf U)|
  \le \sup_U |\Fib_U(a)|.
\]
In particular, if L1 proves \(\sup_U|\Fib_U(k+\sigma)|\le n^B\), then
\[
  |\Lam(C^{\equiv\mu},1-\rho-\sigma/n)|\le n^B
\]
for every ordinary column-distance protocol arity \(\mu\).
\end{lemma}

\begin{proof}
Fix a total order on \(D\).  A listed interleaved codeword
\((P_1,\ldots,P_\mu)\in C^{\equiv\mu}\) at column distance at most \(\delta\) from
\(\mathbf U\) agrees with \(\mathbf U\) on at least \(a\) common columns.  Choose
the first \(a\)-subset \(S\) of those agreeing columns.  For this \(S\), every
row \(U_i|_S\) is interpolated by the degree-\(<k\) polynomial \(P_i\), hence
\(S\in\Fib^{(\mu)}_{\mathbf U}(a)\).

For fixed \(S\), the tuple \((P_1,\ldots,P_\mu)\) is uniquely determined,
because each \(P_i\) has degree less than \(k\) and is fixed by its values on
\(|S|\ge k\) points.  Thus the canonical-support map from listed interleaved
codewords to common feasible supports is injective.  The second inequality is
immediate because every common feasible support is feasible for each individual
row.  Taking suprema gives the final claim.
\end{proof}

\begin{remark}[impact]
For ordinary column-distance RS interleaving, L2 should be demoted from a
separate product-exponent problem to a corollary of L1.  The hard list theorem
remains the base arbitrary-word locator local limit.
\end{remark}

\section{New Bridge II: List plus MCA Gives Line Decoding}
\label{sec:list-plus-mca-gives-line-decoding}

\textbf{Status: \textsc{New-local} bridge.}
The next proposition turns a line-decoding proof obligation into a finite-constant
budget once L1 and M1 are available.  It uses finite slopes and the support-wise
MCA normalization used in this packet.

\begin{proposition}[list/MCA bridge to line decoding]
\label{prop:list-mca-to-line-decoding-bridge}
Let \(C=\RS[\F,D,k]\), let \(|D|=n\), and let
\(a=\lceil(1-\delta)n\rceil>k\).  Suppose that every allowed received line
\(f+\gamma g\) has at most \(M\) support-wise MCA-bad finite slopes at radius
\(\delta\), equivalently \(\epsmca(C,\delta)\le M/|\F|\) for the chosen finite
slope sampler.  Suppose also that the two-row interleaved list size satisfies
\[
  |\Lam(C^{\equiv2},\delta,W)|\le L
  \qquad\text{for every }W\in(\F^2)^D.
\]
Then \(C\) is line-decodable with parameters
\[
  (\delta,\ M+nL+1,\ n+1)
\]
in the following sense: if an allowed received line \(f+\gamma g\) has more than
\(M+nL\) finite slopes \(\gamma\) that are \(\delta\)-close to selected codewords
\(U(\gamma)\in C\), then there are codewords \(c_1,c_2\in C\) such that
\[
  U(\gamma)=c_1+\gamma c_2
\]
for at least \(n+1\) of those slopes.
\end{proposition}

\begin{proof}
For each close slope \(\gamma\), choose a support \(S_\gamma\subseteq D\) of size
\(a\) on which \(f+\gamma g\) agrees with the selected codeword \(U(\gamma)\).
Discard the at most \(M\) support-wise MCA-bad slopes.  For every remaining
slope, support-wise containment gives codewords \(c_{1,\gamma},c_{2,\gamma}\in C\)
such that
\[
  c_{1,\gamma}|_{S_\gamma}=f|_{S_\gamma},
  \qquad
  c_{2,\gamma}|_{S_\gamma}=g|_{S_\gamma}.
\]
Thus \((c_{1,\gamma},c_{2,\gamma})\) is a two-row interleaved codeword
\(\delta\)-close to the received two-row word \((f,g)\), using the common support
\(S_\gamma\).  There are at most \(L\) possible pairs.

If the original line has more than \(M+nL\) close slopes, then after discarding
bad slopes more than \(nL\) slopes remain.  By pigeonhole, one pair
\((c_1,c_2)\) occurs for at least \(n+1\) slopes.  For each such slope, the
codewords \(U(\gamma)\) and \(c_1+\gamma c_2\) both agree with \(f+\gamma g\) on
\(S_\gamma\), hence agree with each other on \(a>k\) positions.  Since both have
degree less than \(k\), they are identical codewords.  This proves the stated
line-decoding conclusion.
\end{proof}

\begin{remark}[finite-constant consequence]
If M1 gives \(M=n^{1+o(1)}\) and L1 plus
Lemma~\ref{lem:common-support-interleaved-fiber-injection} gives \(L=n^B\), then
Proposition~\ref{prop:list-mca-to-line-decoding-bridge} gives
\[
  M+nL+1\le n^{B+1+o(1)}.
\]
For \(n\le2^{44}\) and \(|\F|<2^{256}\), a raw list numerator \(n^B\) is
compatible with a \(2^{-128}\) field budget only for about \(B\lesssim2.9\), and
the line-decoding bridge is compatible only for about \(B\lesssim1.9\).  Thus the
bridge is useful, but it leaves a real finite-constant target.
\end{remark}

\section{Updated Positive-Theorem Package}
\label{sec:updated-positive-theorem-package}

The clean theorem to aim for is the following finite-reserve package.  Let
\[
  C=\RS[\F_{\rm line},D,k],\qquad n=|D|,
  \qquad a=k+\sigma,
  \qquad \delta=1-a/n=1-\rho-\eta.
\]
Assume
\[
  \sigma\log_2\qgen
  \ge (1+\varepsilon)\log_2\binom n{k+\sigma},
  \qquad
  \sigma\ge C_{\rho,B,\varepsilon}\frac n{\log n},
\]
the quotient profile \(\QH(\eta)\) is empty or explicitly budgeted, \(\eta\) lies
above all proved failure floors including Paper D's universal-cap floor, and the
MCA or line experiment is proved over \(\qline\), not merely over \(\qchal\).

\begin{conjecture}[corrected finite reserve package]
\label{conj:corrected-finite-reserve-package}
Under the preceding hypotheses, the desired conclusions are
\[
  |\Lam(C^{\equiv\mu},\delta)|\le n^B
\]
for all ordinary column-distance protocol arities \(\mu\), by
Lemma~\ref{lem:common-support-interleaved-fiber-injection}; and
\[
  \epsmca(C,\delta)
  \le
  \frac{
    n^{1+o(1)}
    +2^{(\beta(\rho)/\Htwo(\rho))\QH(\eta)(1+o(1))}
  }{\qline},
\]
or the same bound over \(\qgen\) plus a proved extension-line lift.  If the
consumed primitive is line decoding, then
Proposition~\ref{prop:list-mca-to-line-decoding-bridge} should give
\[
  C\text{ is }(\delta,\ M+nL+1,\ n+1)\text{ line-decodable},
\]
with \(M\) and \(L\) supplied by M1 and L1.
\end{conjecture}

This package would turn the current negative cap into a sharp threshold program:
below the cap one cannot hope for \(2^{-128}\) MCA soundness, while above the
corrected reserve one would have finite numerators that can be checked against
\(\qline\).

\section{Immediate Experimental Additions}
\label{sec:immediate-experimental-additions}

\begin{enumerate}
\item Add \texttt{experimental/interleaved\_fiber\_injection.md}, containing
Lemma~\ref{lem:common-support-interleaved-fiber-injection}.  Update certificate
logic so an L1 locator-fiber theorem gives \(L_\mu=n^B\), not \(n^{\mu B}\), for
ordinary column-distance RS interleaving.

\item Add \texttt{experimental/list\_mca\_to\_linedecoding.md}, containing
Proposition~\ref{prop:list-mca-to-line-decoding-bridge}.  This gives a concrete
M2 bridge and turns line decoding into a finite-constant ledger problem.

\item Build a certificate scanner reporting generated-field entropy margin,
quotient profile after \(k=\rho n-r\), Paper D cap floor, base and interleaved
list numerators, MCA numerator, line-decoding numerator, field separation among
\(\qgen\), \(\qline\), and \(\qchal\), extension status, and exact failure audit.

\item Keep CA/list statements that use Crites--Stewart marked
\textsc{Conditional} until audited.  Paper D v5's MCA cap should instead be
tracked as \textsc{Proved} when the divisor, binomial, and subfield hypotheses
are verified.
\end{enumerate}

\section{Fresh Non-Claims}
\label{sec:fresh-nonclaims}

This revision does not claim an exact settlement of the smooth-case Proximity
Prize threshold.  It claims that the shape is now rigid:
\[
  \boxed{\text{No unslacked capacity theorem; prove a reserve theorem at }
  1-\rho-\eta.}
\]
The theorem-backed negative side gives a draft universal MCA cap at constant gap
\(2^{-9}\) for \(\rho=1/2,1/4,1/8\) and \(2^{-10}\) for \(\rho=1/16\).  The
positive side should be pursued as
\[
  \text{L1 locator local limit}
  +\text{M1 residue-line packing}
  +\text{extension-line theorem}
  +\text{protocol-ledger rewrite}.
\]
The interleaved-list problem is less fundamental than it looked in the old file:
for column-distance interleaving, a uniform locator-fiber theorem immediately
bounds the interleaved list at the same exponent.

\section{Retained Previous Experiments Material}
\label{sec:retained-previous-experiments-material}

The remaining sections are retained from the previous \texttt{experiments.tex}
whenever they still provide useful provenance, notation, normal forms, finite
certificate gates, or non-claims.  The new front matter above should be read as
the current threshold-level guide; the older material below remains a toolbox.

\section{Status, Provenance, and What Counts as New}

This file is the place where we collect new things, but it must separate new
repository contributions from imported constructions.  The status labels used
below are:
\[
\begin{array}{ll}
\textsc{Imported} & \text{an external theorem, proof pattern, or conversion;}\\
\textsc{Wrapper} & \text{a repository repackaging, dictionary, or pigeonhole consequence;}\\
\textsc{New-local} & \text{a restricted algebraic theorem proved in this note;}\\
\textsc{Target} & \text{a precise proof obligation not yet promoted;}\\
\textsc{Heuristic} & \text{evidence or a search direction, not a theorem.}
\end{array}
\]

\paragraph{Hard attribution rule.}
Any locator-polynomial argument derived from the subgroup construction of
Ben-Sasson--Carmon--Hab\"ock--Kopparty--Saraf must cite \cite[Thm.~7.1]{BCHKS25}
where the construction is stated or proved.  Any admissible subgroup
instantiation must also cite \cite[Thm.~1.13]{BCHKS25}.  The citation belongs
at the theorem statement and again at the proof entry point; a bibliography
entry or acknowledgement alone is not enough.

\paragraph{What is new in this note.}
The following items are the repository-side contributions collected here.
\begin{itemize}
\item A smooth-quotient notation dictionary translating BCHKS subgroup
  notation into the RS-MCA notation \(Q=D^{\efib}\) and the dyadic/FFT-domain
  ledgers used in Papers A--D.
\item A list-fiber pigeonhole wrapper: once a locator construction assigns
  many subset certificates to a smaller slope set, one heavy slope has many
  locator certificates, and, after a distinctness check, many nearby
  codewords.
\item A slack-two and subfield-pigeonhole proof target for the older
  Crites--Stewart/CA route: the imported locator construction should be run at
  the augmented-code rung needed by that conversion, and slopes should be
  pigeonholed over the generated/base field when the construction really takes
  values there.
\item A \textsc{New-local} divisibility gate in the restricted Paper B window
  \(B=\F_p\), \(F=\F_{p^2}\), \(D=\F_p\), \(t=\sigma=2\), \(j=3\).  This is
  independent of the BCHKS priority issue and is the main theorem proved in
  Section~\ref{sec:divisibility-gate}.
\end{itemize}

\paragraph{What is not new.}
The polynomial identity behind the item-(d) locator proof is not new here.  In
BCHKS notation it is the expansion of a vanishing polynomial over a complement
of a chosen subset; in smooth-quotient notation it is the same expansion of
\(\prod_{b\in A}(X^{\efib}-b)\), up to choosing a subset or its complement.
The use of \(\ell\) for the chosen subset size also follows the external
presentation closely enough that it should either be cited or renamed to
\(\ellpg\).

\paragraph{Source map.}
The main papers remain the authoritative documents.
\begin{itemize}
\item Paper A, \texttt{tex/RS\_disproof\_v3.tex}, refutes the old no-slack
  smooth-domain support-wise line-MCA statement.  If it uses the locator
  construction, it should cite BCHKS and call the repository contribution a
  smooth-domain/list-fiber corollary, not a new locator proof.
\item Paper B, \texttt{tex/slackMCA\_v3.tex}, is the main target for the local
  resonance material in Sections~\ref{sec:resonance-setup}--\ref{sec:scanner}.
\item Paper C, \texttt{tex/snarks\_v4.tex}, records protocol ledgers.  Nothing
  in this note may be consumed by a protocol without a separate conversion
  theorem.
\item Paper D, \texttt{tex/cs25\_cap\_v4.tex}, gives the universal cap via an
  imported list-to-agreement route.  Its locator-side input should cite BCHKS
  explicitly if it uses the quotient-locator construction.
\end{itemize}

\section{BCHKS Locator Construction in Smooth-Quotient Notation}
\label{sec:bchks}

\textbf{Status: \textsc{Imported} plus \textsc{Wrapper}.}
This section is a provenance-safe restatement and dictionary.  It is included
so later main-paper edits have a clean place to copy from.

\subsection{Imported theorem/proof pattern}

BCHKS prove a general multiplicative-subgroup construction
\cite[Thm.~7.1]{BCHKS25}.  In a compact form, the imported statement is as
follows.

\begin{theorem}[BCHKS quotient-locator construction, imported]
\label{thm:bchks-import}
Let \(G\subseteq H\subseteq \F_q^*\) be multiplicative subgroups and let
\(\Phi:H\to G\) be \(x\mapsto x^c\), where \(c=|H|/|G|\).  Let
\(E\subseteq G\), and let \(\ellpg,a\geq1\) satisfy
\[
  |E^{(+\ellpg)}|
  =\#\{e_1+\cdots+e_{\ellpg}:e_i\in E\text{ distinct}\}\geq a.
\]
Let \(D=\Phi^{-1}(E)\), \(n=|D|\), and
\[
  C=\RS[\F_q,D,n-(\ellpg+2)c].
\]
Then, for \(\gamma=\ellpg c/n\), there are functions
\(f,g:D\to\F_q\) such that
\[
  \#\{z\in\F_q: \dist(f+zg,C)\leq \gamma\}\geq a,
\]
while
\[
  \dist([f,g],C^2)\geq (\ellpg+1)c/n.
\]
\end{theorem}

\begin{proof}[Reference]
This is exactly the imported construction of \cite[Thm.~7.1]{BCHKS25}; the
proof is not reproduced as a repository contribution.  BCHKS define the line
endpoints by the monomials \(x^{n-\ellpg c}\) and
\(x^{n-(\ellpg+1)c}\), then expand a vanishing polynomial over the complement
of a chosen \(\ellpg\)-subset.  The coefficient of
\(X^{n-(\ellpg+1)c}\) is an affine transform of a restricted sum from
\(E^{(+\ellpg)}\), giving the exceptional slope values.
\end{proof}

The BCHKS Theorem~1.13 is the admissible-subgroup specialization of this
construction: take \(E=G\), take \(H=D\), set \(c=|D|/|G|\), and use
\(\ellpg=|G|/2\).  Any RS-MCA statement corresponding to that specialization
should cite \cite[Thm.~1.13]{BCHKS25} directly.

\subsection{Notation dictionary}

\begin{center}
\begin{tabular}{p{0.28\linewidth}p{0.30\linewidth}p{0.32\linewidth}}
\textbf{BCHKS notation} & \textbf{RS-MCA notation} & \textbf{Meaning}\\
\hline
\(H\) & evaluation domain \(D\) & Large multiplicative domain on which codewords are evaluated.\\
\(G\) & quotient \(Q=D^{\efib}\) & Image of the power map; one point in \(Q\) has \(\efib\) preimages in \(D\).\\
\(c=|H|/|G|\) & \(\efib=n/|Q|\) & Fiber size of the quotient map.  Avoid calling this \(a\), because BCHKS also uses \(a\) for the number of exceptional slopes.\\
\(\Phi(x)=x^c\) & \(x\mapsto x^{\efib}\) & Quotient/fiber map.\\
\(E\subseteq G\) & active quotient subset \(E\subseteq Q\) & Quotient-level support used by the locator.\\
\(E^{(+\ellpg)}\) & restricted first-coefficient/slope set & Sums of distinct chosen quotient elements; slopes are affine transforms of these sums.\\
\(\ellpg\) & \(\ellpg\) or renamed subset size & Number of chosen quotient elements.  Use \(\ellpg\) in RS-MCA text to avoid collision with other \(\ell\)'s.\\
\(f=x^{n-\ellpg c}\), \(g=x^{n-(\ellpg+1)c}\) & locator line endpoints & Same monomial line after quotient notation is substituted.\\
\(P_S(X)=\prod_{\alpha\in S'}(X^c-\alpha)\) & \(\prod_{b\in A}(X^{\efib}-b)\), possibly after complementing \(A\) & The locator polynomial.\\
exceptional \(z\)'s & bad slopes / heavy list fibers & Slope parameters for which a line word is close to the RS code.\\
\end{tabular}
\end{center}

\subsection{The shared locator identity}

The following identity is written down because it is the exact place where
priority can be lost.  The identity is the BCHKS identity in RS-MCA notation.

\begin{lemma}[Locator expansion, imported identity]
\label{lem:locator-imported}
Let \(D\) map \(\efib\)-to-one onto a quotient set \(E\) by
\(x\mapsto x^{\efib}\), and let \(A\subseteq E\) have size \(\ellpg\).  Put
\(A^c=E\setminus A\) and define
\[
  P_A(X)=\prod_{\alpha\in A^c}(X^{\efib}-\alpha).
\]
Then
\[
  P_A(X)
  =X^{n-\ellpg\efib}
    -e_1(A^c)X^{n-(\ellpg+1)\efib}
    +\text{terms of degree at most }n-(\ellpg+2)\efib.
\]
Equivalently, since \(e_1(A^c)=e_1(E)-e_1(A)\), the coefficient of the next
monomial is an affine transform of the restricted sum \(e_1(A)\).
\end{lemma}

\begin{proof}[Reference]
This is the elementary expansion used inside \cite[Thm.~7.1]{BCHKS25}.  We
record it only to fix notation.  It should not be cited as a new lemma of the
repository unless the BCHKS provenance is present at the statement.
\end{proof}

\subsection{List-fiber pigeonhole wrapper}

The next proposition is the repository-side wrapper around the imported
identity.  It is intentionally stated in certificate form, because distinctness
of the nearby codewords is a separate check in concrete instantiations.

\begin{proposition}[List-fiber pigeonhole wrapper]
\label{prop:list-fiber-wrapper}
Let \(\mathcal A\) be a family of \(\ellpg\)-subsets of a quotient set \(E\).
Suppose the BCHKS locator construction assigns to each \(A\in\mathcal A\) a
slope \(z(A)\) lying in a finite set \(S\), and a locator certificate showing
that the corresponding word \(f+z(A)g\) is within radius \(\gamma\) of a
locator codeword \(P_A\in C\).  Then some slope \(z_0\in S\) has at least
\[
  \frac{|\mathcal A|}{|S|}
\]
locator certificates above it.  If the codewords \(P_A\) attached to the
certificates over \(z_0\) are distinct, then the Hamming ball of radius
\(\gamma\) around \(f+z_0g\) contains at least \(|\mathcal A|/|S|\) codewords
of \(C\).
\end{proposition}

\begin{proof}
The first assertion is the pigeonhole principle applied to the map
\(A\mapsto z(A)\).  The second assertion is only the definition of list size,
after the distinctness check removes possible duplicate locator codewords.
The construction that supplies the certificates is imported from
Theorem~\ref{thm:bchks-import}.
\end{proof}

\begin{remark}[What the wrapper adds]
The new part here is not the locator polynomial.  The wrapper isolates the
extra bookkeeping needed by RS-MCA: which slope set pays the pigeonhole bill,
whether the slope set is the ambient field \(\F\) or a generated/base subfield
\(B\), and whether locator certificates remain distinct after specialization.
Those are the checks needed before using the construction in Papers A, B, or D.
\end{remark}

\subsection{Slack-two and subfield target for the older CS25 route}

\textbf{Status: \textsc{Target}.}
The older CS25/CA route needs a locator-heavy list fiber in the augmented code
rung \(C^+=\RS[\F,D,k+1]\), because that is the object fed into the
Crites--Stewart list-to-agreement conversion.  Paper D v5 no longer uses this
as the load-bearing MCA-cap route, but the target remains useful for CA/list
comparison audits.  The proof obligation should be stated as follows.

\begin{proposition}[Slack-two/subfield locator target]
\label{prop:slack-two-target}
Let \(D\subseteq B\subseteq\F\), and suppose a BCHKS-type quotient locator
construction over \(D\) produces slopes in \(B\), not merely in \(\F\).  If the
one-rung/augmented-code version produces distinct locator codewords in
\(C^+=\RS[\F,D,k+1]\) for a family \(\mathcal A\) of certificates, then there
exists a slope \(z_0\in B\) with list size at least
\[
  \frac{|\mathcal A|}{|B|}
\]
inside the appropriate Hamming ball for \(C^+\).
\end{proposition}

\begin{proof}[Bookkeeping only]
This is Proposition~\ref{prop:list-fiber-wrapper} with \(S=B\).  The nontrivial
work is external or elsewhere: one must verify the BCHKS locator hypotheses at
the augmented-code rung, verify that the slope coefficients lie in \(B\), and
verify distinctness of the produced codewords.  Until those checks are present,
this proposition is a target, not a promoted Paper D theorem.
\end{proof}

\subsection{Pasteable main-paper wording}

\paragraph{Paper A wording.}
The polynomial locator identity used here is the quotient-fiber form of the
construction of Ben-Sasson--Carmon--Hab\"ock--Kopparty--Saraf
\cite[Thm.~7.1]{BCHKS25}; their Theorem~1.13 is the admissible subgroup
instantiation.  Thus the locator proof itself is not a new contribution of
this manuscript.  The present use is the smooth-domain specialization, the
field/divisor bookkeeping, and the list-fiber pigeonhole packaging.

\paragraph{Paper D wording.}
The locator-side input is the BCHKS quotient-locator construction
\cite[Thm.~7.1; cf.~Thm.~1.13]{BCHKS25}, written in the smooth-domain quotient
notation \(Q=D^{\efib}\).  The additional bookkeeping needed here is the
slack-two augmented-code rung and, when available, subfield pigeonholing over
\(B\) rather than over the ambient extension field \(\F\).

\section{Triage of the AI-Generated Four-Item Packet}
\label{sec:packet-triage}

\textbf{Status: \textsc{Audit}.}
The four advertised items should be positioned as follows before any promotion
into a main paper.  The original AI packet used the labels (a)--(d) without
making them self-contained in this note, so we decode them here.

\begin{center}
\begin{tabular}{p{0.10\linewidth}p{0.28\linewidth}p{0.50\linewidth}}
\textbf{Item} & \textbf{Short name} & \textbf{Meaning in this ledger}\\
\hline
(a) & Weak-slack positive regime & A positive-looking proximity/list or
MCA-adjacent statement that only reaches slack on the order of
\(1/\sqrt n\).  It is useful as a comparison theorem, but it is far from the
constant/reserve bookkeeping needed for the corrected RS-MCA target.\\
(b) & Finite Fermat-prime packet & A special finite-domain instantiation over
only three Fermat-prime cases, in the current experimental ledger represented
by the \(p\in\{17,257,65537\}\) style checks.  It is evidence or a worked
computation, not a uniform theorem.\\
(c) & Exponential-field construction & A large-field variant whose proof
requires the ambient field to be exponential in the domain/parameter size.
This can clarify the mechanism but does not match the polynomial-field or
challenge-field regimes needed by the Proximity Prize ledger.\\
(d) & BCHKS quotient-locator packet & The interesting locator/proximity-gap
construction: choose a quotient of a multiplicative subgroup, use
\(\ellpg\)-subsets and restricted sums to create many exceptional slopes or
nearby locator certificates.  This is the BCHKS construction of
\cite[Thm.~7.1]{BCHKS25}, with the admissible subgroup instance in
\cite[Thm.~1.13]{BCHKS25}; the repository contribution is only the
smooth-domain dictionary, list-fiber wrapper, and later slack-two/subfield
bookkeeping.\\
\end{tabular}
\end{center}

Item (d) should not be confused with the \textsc{New-local} divisibility-gate
theorem in Section~\ref{sec:divisibility-gate}.  The former is an imported
locator/proximity construction; the latter is an independent restricted Paper
B resonance calculation.

\begin{center}
\begin{tabular}{p{0.10\linewidth}p{0.34\linewidth}p{0.44\linewidth}}
\textbf{Item} & \textbf{Limitation} & \textbf{Repository handling}\\
\hline
(a) & Works only at slack on the order of \(1/\sqrt n\). & Use as a weak-slack comparison point only.  It does not clear the corrected-reserve MCA objective.\\
(b) & Applies only to three Fermat-prime cases. & Treat as finite-prime evidence or a worked computation, not as a uniform smooth-domain theorem.\\
(c) & Requires exponential field size. & Treat as a large-field sanity check.  It is not directly usable for polynomial-field or prize-ledger parameters.\\
(d) & Most relevant, but the locator proof is essentially BCHKS. & Import, cite, and adapt from \cite[Thm.~7.1 and Thm.~1.13]{BCHKS25}.  The repository contribution is only the smooth-domain/list-fiber/slack-two packaging after exact hypotheses are checked.\\
\end{tabular}
\end{center}

\paragraph{Promotion checklist.}
Before any statement using item (d) moves to Papers A--D, check the following.
\begin{itemize}
\item Does the statement cite BCHKS at the theorem and proof entry point?
\item Is the object a proximity-gap statement, a list statement, CA, MCA,
  line-decoding, or a protocol-facing theorem?
\item Which field pays the slope/list pigeonhole bill: \(B\), \(\F\), or
  another generated field?
\item Are the locator codewords distinct after specialization?
\item Is the claimed slack in the same normalization as the consumed theorem?
\item Are items (a), (b), and (c) kept out of asymptotic/protocol claims unless
  their limitations are explicitly part of the theorem statement?
\end{itemize}

\section{Restricted Resonance Setup for Paper B}
\label{sec:resonance-setup}

\textbf{Status: \textsc{New-local} for the divisibility gate; otherwise local
normal-form bookkeeping.}
The remaining sections preserve and reorganize the Cycle 14--18 algebraic
material.  This branch is independent of the BCHKS priority correction.

The common restricted setting is
\[
  B=\F_p,\qquad F=\F_{p^2},\qquad D=\F_p,\qquad
  t=\sigma=2,\qquad j=3,
\]
and we work off the tangent/global endpoint
\[
  R_0=\{\, [W]_E\wedge [\Bnum]_E=0\,\}.
\]
The field ledgers remain distinct:
\[
  \qgen=p,\qquad \qline=p^2.
\]
The regime is sub-reserve, since \(\eta=\sigma/n=2/n\).  Thus none of the
following proves a corrected-reserve list, CA, MCA, line-decoding, SNARK, or
protocol theorem.

The base/generated field is \(B=\F_p\).  The line field is the quadratic
extension \(F=\F_{p^2}=B(\alpha)\), with basis \(\{1,\alpha\}\) over \(B\) and
\(\alpha^2=\nu\in B\).  The domain is \(D=\F_p\), so \(n=|D|=p\).

A degree-two residue-line datum consists of
\[
  E\in F[X],\qquad \deg E=2,
  \qquad E(x)\ne0\ \text{for all }x\in D,
\]
and a numerator
\[
  \Bnum\in F[X],\qquad \deg \Bnum<2.
\]
We write
\[
  A=F[X]/E,
  \qquad [P]_E\in A,
  \qquad b=[\Bnum]_E.
\]
The anchor word on \(D\) is interpolated by a polynomial \(W\).  The operation
\(u\wedge v\) means the determinant of the coordinate vectors of \(u,v\in A\)
in a chosen \(F\)-basis of the two-dimensional \(F\)-space \(A\).  Off \(R_0\),
the pair \(\{[W]_E,b\}\) is such a basis.

The co-support has size \(j=3\).  If \(T\subset D\) is the three-point
co-support, define
\[
  \tau_1=e_1(T),\qquad \tau_2=e_2(T),\qquad \tau_3=e_3(T).
\]
All valid co-supports have \(\tau_i\in B\).  The quantity \(\Ctwo\) denotes the
number of distinct slopes produced by the relevant landing branch after
restricting to valid split triples \(T\subset\F_p\).  Bounding landing points
therefore bounds \(\Ctwo\), but a sharp slope theorem may require more because
many landing points can share a slope.

\section{Cycle 14 Normal Form}
\label{sec:cycle14}

The Cycle 14 audit reduces the \(t=2,j=3\) landing equation to two affine forms
in the elementary co-support parameters
\[
  \tau=(\tauone,\tautwo,\tauthree).
\]

\begin{lemma}[Affine wedge normal form]
\label{lem:affine-wedge}
Assume \(\{[W]_E,b\}\) is an \(F\)-basis of \(A\).  Then the landing equation can
be written in the form
\[
  \iota(\tau)=A_0(\tauone,\tautwo)-\tauthree [W]_E,
  \qquad
  \mu(\tau)=B_0(\tauone,\tautwo)-\tauthree b,
\]
where
\[
  A_0=p_1[W]_E+p_2 b,
  \qquad
  B_0=q_1[W]_E+q_2 b,
\]
and \(p_i,q_i\in F\) are affine-linear functions of
\((\tauone,\tautwo)\).  If the basis is normalized so \([W]_E\wedge b=1\), then
\[
  \Delta=\iota\wedge\mu
    =(p_1-\tauthree)(q_2-\tauthree)-p_2q_1.
\]
\end{lemma}

\begin{proof}
The Cycle 14 quotient calculation, which is a concrete instance of Paper B's
residue-line normal form, gives the displayed affine forms.  Expanding in the
ordered basis \(\{[W]_E,b\}\), the coordinate vector of \(\iota\) is
\((p_1-\tauthree,p_2)\), while the coordinate vector of \(\mu\) is
\((q_1,q_2-\tauthree)\).  Their wedge determinant is therefore
\[
  (p_1-\tauthree)(q_2-\tauthree)-p_2q_1.
\]
\end{proof}

\section{Cycle 18 Component Split}
\label{sec:cycle18}

Choose a basis \(F=B+\alpha B\) with \(\alpha^2=\nu\in B\), and write
\[
  p_1=p_{10}+\alpha p_{11},\quad
  p_2=p_{20}+\alpha p_{21},\quad
  q_1=q_{10}+\alpha q_{11},\quad
  q_2=q_{20}+\alpha q_{21}.
\]
Split
\[
  \Delta=\Deltazero+\alpha\Deltaone,
  \qquad
  \Deltazero,\Deltaone\in B[\tauone,\tautwo,\tauthree].
\]

\begin{theorem}[Monic quadratic and linear alpha component]
\label{thm:component-split}
In the setting above,
\[
\begin{aligned}
\Deltazero={}&
  \tauthree^2-(p_{10}+q_{20})\tauthree
  +p_{10}q_{20}+\nu p_{11}q_{21}
  -p_{20}q_{10}-\nu p_{21}q_{11},\\
\Deltaone={}&
  -(p_{11}+q_{21})\tauthree
  +p_{10}q_{21}+p_{11}q_{20}-p_{20}q_{11}-p_{21}q_{10}.
\end{aligned}
\]
Consequently \(\Deltazero\) is monic quadratic in \(\tauthree\), while
\(\Deltaone\) has degree at most one in \(\tauthree\).
\end{theorem}

\begin{proof}
Substitute the coordinate expressions for \(p_i,q_i\) into the identity of
Lemma~\ref{lem:affine-wedge} and use \(\alpha^2=\nu\).  Collecting the
coefficient of \(1\) gives the displayed formula for \(\Deltazero\), and
collecting the coefficient of \(\alpha\) gives the displayed formula for
\(\Deltaone\).  The coefficient of \(\tauthree^2\) in \(\Deltazero\) is \(1\),
and the \(\alpha\)-component has no \(\tauthree^2\)-term.
\end{proof}

\begin{corollary}[Graph reduction outside the zero-alpha branch]
\label{cor:graph}
If \(\Deltaone\) is not the zero polynomial, write
\[
  \Deltaone=s(\tauone,\tautwo)\tauthree+h(\tauone,\tautwo).
\]
On the open locus \(s\ne0\), the common-zero condition
\(\Deltazero=\Deltaone=0\) is supported on the graph
\[
  \tauthree=-h/s.
\]
\end{corollary}

\begin{proof}
This is immediate from solving the linear equation \(\Deltaone=0\) for
\(\tauthree\) on \(s\ne0\).
\end{proof}

\section{New Divisibility-Gate Theorem}
\label{sec:divisibility-gate}

The next theorem is the main \textsc{New-local} theorem extracted from the
Cycle 18 reduction.  It is stronger than merely saying that the branch is a
graph: it shows that any two-dimensional family must satisfy an explicit
divisibility identity.

\begin{theorem}[Divisibility gate for the graph branch]
\label{thm:divisibility-gate}
Assume the restricted setting above and assume \(\Deltaone\ne0\).  Write
\[
  \Deltaone=s(\tauone,\tautwo)\tauthree+h(\tauone,\tautwo),
\]
and define the denominator-cleared graph polynomial
\[
  G(\tauone,\tautwo)
  =s(\tauone,\tautwo)^2
    \Deltazero\bigl(\tauone,\tautwo,-h(\tauone,\tautwo)/s(\tauone,\tautwo)\bigr).
\]
Then \(\deg G\le 4\).  If \(G\) is not the zero polynomial, then
\[
  \#\{\,\tau\in B^3:\Deltazero(\tau)=\Deltaone(\tau)=0\,\}=O(p).
\]
Consequently the corresponding number of distinct slopes satisfies
\[
  \Ctwo=O(p)=O(n).
\]
Therefore a growing-prime family with
\[
  \Ctwo=\Theta(p^2)=\Theta(\qline)
\]
in the nonzero-\(\Deltaone\) branch can occur only if \(G\equiv0\).
\end{theorem}

\begin{proof}
Because the \(p_i,q_i\) are affine-linear in \((\tauone,\tautwo)\), the
coefficient \(s\) has degree at most \(1\), and \(h\) has degree at most \(2\).
Since \(\Deltazero\) is monic quadratic in \(\tauthree\), clearing the
denominators after substituting \(\tauthree=-h/s\) gives a polynomial \(G\) of
degree at most \(4\).

On \(s\ne0\), the equation \(\Deltaone=0\) forces \(\tauthree=-h/s\).  Hence
every common zero on this open locus gives a base pair
\((\tauone,\tautwo)\) with \(G(\tauone,\tautwo)=0\).  If \(G\) is nonzero,
Schwartz--Zippel gives at most \(4p\) such base pairs, and each base pair gives
at most one value of \(\tauthree\).

It remains to count the exceptional locus \(s=0\).  There \(\Deltaone=0\) also
forces \(h=0\).  If \(s\) is a nonzero affine form, this is contained in an
affine line; if \(s\equiv0\), then \(h\) is a nonzero polynomial of degree at
most \(2\).  In either case there are \(O(p)\) base pairs.  For each such pair,
\(\Deltazero\) is a monic quadratic in \(\tauthree\), so it gives at most two
values of \(\tauthree\).  The total common-zero count is \(O(p)\).

Finally, the number of distinct slopes is bounded by the number of landing
points \(\tau\).  Thus \(\Ctwo=O(p)\) whenever \(G\not\equiv0\), and any
\(\Theta(p^2)\) family in this branch must satisfy \(G\equiv0\).
\end{proof}

\begin{remark}
The identity \(G\equiv0\) says that, after clearing the \(s\)-denominator,
\(\Deltazero\) is divisible by the linear equation \(\Deltaone\) on the graph
branch.  This is the first exact algebraic gate that a large slope family in
this restricted window must pass.
\end{remark}

\section{Nearby Proven Gates}
\label{sec:nearby-gates}

The same restricted lane contains two other useful gates.  They are included
here because they explain where Theorem~\ref{thm:divisibility-gate} fits.

\begin{proposition}[Base-component complete-intersection gate]
\label{prop:complete-intersection}
Suppose \(\Deltazero,\Deltaone\in B[\tauone,\tautwo,\tauthree]\) have no common
surface component.  Then
\[
  \#V(\Deltazero,\Deltaone)(B)=O(p),
\]
and hence this branch contributes \(\Ctwo=O(p)\) slopes.
\end{proposition}

\begin{proof}
The two components have bounded degree.  If they share no surface component,
their common zero set in \(\mathbb A^3_B\) is a curve of bounded degree.  A
bounded-degree affine curve over \(\F_p\) has \(O(p)\) rational points, and
slopes are bounded by landing points.
\end{proof}

\begin{proposition}[Cycle 16 determinant gate]
\label{prop:det-gate}
Let \(Q(z_0,z_1)\) be the Cycle 15--16 determinant consistency polynomial for
the affine system
\[
  L_z(\tau)=\iota(\tau)-z\mu(\tau)=0,
  \qquad z=z_0+\alpha z_1.
\]
If \(Q\) is not identically zero, then in the restricted \(D=\F_p\),
\(t=\sigma=2\), \(j=3\), off-\(R_0\) window,
\[
  \Ctwo\le 4p=O(p)=O(n).
\]
\end{proposition}

\begin{proof}
Cycle 16 gives \(\deg Q\le4\) and shows that every landing slope must satisfy
\(Q(z_0,z_1)=0\).  If \(Q\ne0\), Schwartz--Zippel over \(B=\F_p\) gives at most
\(4p\) possible pairs \((z_0,z_1)\), hence at most \(4p\) slopes.
\end{proof}

\section{Slope Map and Remaining Heuristics}
\label{sec:slope-map}

On the graph branch, the landing equation \(\iota=z\mu\) gives
\[
  p_1-\tauthree=zq_1,
  \qquad
  p_2=z(q_2-\tauthree).
\]
Where \(q_1\ne0\), substituting \(\tauthree=-h/s\) yields
\[
  z(\tauone,\tautwo)=\frac{p_1+h/s}{q_1}.
\]
Equivalently, eliminating \(\tauthree\) recovers the Cycle 14 quadratic
\[
  q_1z^2-(p_1-q_2)z-p_2=0,
\]
so the residual map is controlled by
\[
  (\tauone,\tautwo)
  \longmapsto [\,q_1:(p_1-q_2):p_2\,]\in\mathbb P^2(F).
\]

\begin{heuristic}[Exceptional identity should be rare]
For generic source data, the polynomial \(G\) should not vanish identically.
Thus the graph branch should usually be curve-sized, and the first possible
large-slope families should come from algebraically special data satisfying
the divisibility identity \(G\equiv0\) or from the zero-alpha branch
\(\Deltaone\equiv0\).
\end{heuristic}

\begin{heuristic}[What remains after the gate]
If \(G\equiv0\), the problem is no longer to show that the common-zero set is
small.  It may be surface-sized.  The remaining question is whether the
projective coefficient map has one-dimensional image on every source-valid
surface, giving \(\Ctwo=O(p)\), or whether some growing-prime source-valid
family has two-dimensional image and \(\Ctwo=\Theta(p^2)\).
\end{heuristic}

\section{Next Experimental Checks}
\label{sec:scanner}

A useful scanner should compute, for each source-valid candidate,
\[
  \Deltaone=s\tauthree+h,
  \qquad
  G=s^2\Deltazero(\tauone,\tautwo,-h/s),
\]
then record whether \(G\equiv0\), whether \(\Deltaone\equiv0\), the exceptional
locus \(s=h=0\), the image size of
\[
  [\,q_1:(p_1-q_2):p_2\,],
\]
and the direct slope count over distinct split triples \(T\subset\F_p\).

The certificate should keep at least
\[
\begin{gathered}
  p,\ \qgen,\ \qline,\ E,\ \Bnum,\ \text{seed},\ \text{stratum},\
  R_0\text{-status},\\
  \Deltaone\equiv0,
  \quad G\equiv0,
  \quad \#\{s=h=0\},
  \quad \#\operatorname{image},
  \quad \Ctwo,
  \quad \#\{T\text{ split distinct}\}.
\end{gathered}
\]

\section{June 26 Normal-Form and Dependency Results}
\label{sec:june26}

\textbf{Status: mixed.}
This section records the parts of the June 26 PR batch that survived a
statement-level validity check.  The statements below are useful proof
infrastructure, not prize claims.  In particular, they do not improve the
current Cycle116/119 numerator \(52{,}747{,}567{,}092\) over
\(\F_{17^{32}}\), and they do not prove the corrected Paper B MCA conjecture.

\subsection{Syndrome-pencil normal form for support-wise line incidence}

\textbf{Status: \textsc{New-local}.}
This is the cleanest reusable F1 contribution from the batch.  It turns
support-wise line incidence into a Hankel-pencil locator equation over the
actual line field.

Let \(D=\{x_1,\ldots,x_n\}\subset F\) have distinct points and let
\(C=\RS[F,D,k]\) with redundancy \(r=n-k\).  For each \(x_i\), put
\[
  \lambda_i=\left(\prod_{h\ne i}(x_i-x_h)\right)^{-1}.
\]
For a word \(y:D\to F\), define
\[
  \operatorname{Syn}(y)_m
  =\sum_i \lambda_i x_i^m y(x_i),
  \qquad 0\le m<r.
\]
Then \(\operatorname{Syn}(y)=0\) if and only if \(y\in C\).  Fix
\(T\subset D\) with \(|T|=j\), assume \(0\le j\le r\), and set \(t=r-j\).
Write
\[
  L_T(X)=\prod_{x\in T}(X-x)=\ell_0+\ell_1X+\cdots+\ell_jX^j,
  \qquad
  \boldsymbol\ell_T=(\ell_0,\ldots,\ell_j)^T .
\]
For \(w\in F^r\), define the \(t\times(j+1)\) Hankel window
\[
  H_{t,j}(w)_{a,b}=w_{a+b},
  \qquad 0\le a<t,\quad 0\le b\le j.
\]

\begin{theorem}[Hankel-pencil test]
\label{thm:syndrome-pencil}
Let \(f,g:D\to F\), put \(u=\operatorname{Syn}(f)\) and
\(v=\operatorname{Syn}(g)\), and let \(S=D\setminus T\).  For a finite slope
\(z\in F\), the word \(f+zg\) is explained by a codeword of \(C\) on \(S\) if
and only if
\[
  \bigl(H_{t,j}(u)+zH_{t,j}(v)\bigr)\boldsymbol\ell_T=0.
\]
Moreover, this explanation is support-wise noncontained for the line
\(f+zg\) on \(S\) if and only if, in addition,
\[
  H_{t,j}(v)\boldsymbol\ell_T\ne0.
\]
Consequently the support-wise MCA bad slopes at agreement \(k+t\) are exactly
the slopes for which there is a squarefree \(D\)-split locator \(L_T\) of
degree \(j=r-t\) satisfying the displayed incidence and noncontainment tests.
\end{theorem}

\begin{proof}
Let \(W_T\subset F^r\) be the span of the parity-check columns indexed by
\(T\):
\[
  W_T=\operatorname{span}_F\{\lambda_x(1,x,\ldots,x^{r-1}):x\in T\}.
\]
Equivalently, \(w\in W_T\) means that \(w\) is the syndrome of a word supported
on \(T\).  For every \(w\in W_T\), write
\[
  w_m=\sum_{x\in T}c_xx^m.
\]
Then, for \(0\le a<t\),
\[
  \bigl(H_{t,j}(w)\boldsymbol\ell_T\bigr)_a
  =\sum_{b=0}^j\ell_b w_{a+b}
  =\sum_{x\in T}c_xx^aL_T(x)=0.
\]
Conversely, since \(L_T\) is monic, the recurrence
\(\sum_{b=0}^j\ell_b w_{a+b}=0\) for \(0\le a<r-j\) determines all remaining
coordinates from the first \(j\) coordinates, so its solution space has
dimension at most \(j\).  The \(j\) columns indexed by the distinct points of
\(T\) form a Vandermonde system of rank \(j\), hence the recurrence space is
exactly \(W_T\).

The word \(f+zg\) is explained on \(S\) if and only if there is \(c\in C\) such
that \(f+zg-c\) is supported on \(T\).  Taking syndromes, this is equivalent to
\(u+zv\in W_T\), which is the first displayed equation by the recurrence
criterion.  Under that equation, simultaneous explanation of \(f\) and \(g\) on
the same \(S\) is equivalent to \(u,v\in W_T\).  Since
\(u=(u+zv)-zv\), this is equivalent to \(v\in W_T\), hence to
\(H_{t,j}(v)\boldsymbol\ell_T=0\).  Noncontainment is therefore exactly the
negation of that condition.
\end{proof}

\subsection{Interleaved codegree reduction}

\textbf{Status: \textsc{New-local} reduction; saving remains conditional.}
This result is useful for the L2 interleaved-list ledger because it replaces
the naive Cartesian product by a recursion through higher-agreement tails.

Let \(C=\RS[F,D,k]\).  For a word \(V:D\to F\) and codeword \(c\in C\), write
\[
  A_V(c)=\{x\in D:c(x)=V(x)\},
  \qquad
  \operatorname{Fib}_V(s)=\{c\in C:|A_V(c)|\ge s\},
\]
and \(M_V(s)=|\operatorname{Fib}_V(s)|\).  For a \(\mu\)-row word
\(U=(U_1,\ldots,U_\mu)\), define
\[
  \Lambda_\mu^{(s)}(U)
  =
  \#\{(c_1,\ldots,c_\mu)\in C^\mu:
  |\cap_i A_{U_i}(c_i)|\ge s\}.
\]

\begin{theorem}[two-regime codegree bound]
\label{thm:l2-codegree}
For every agreement threshold \(a\),
\[
  \Lambda_2^{(a)}(U_1,U_2)
  \le
  M_{U_2}(a)+M_{U_2}(2a-k)\,M_{U_1}(a).
\]
More generally,
\[
  \Lambda_\mu^{(a)}(U_1,\ldots,U_\mu)
  \le
  \Lambda_{\mu-1}^{(a)}(U_2,\ldots,U_\mu)
  +
  \Lambda_{\mu-1}^{(2a-k)}(U_2,\ldots,U_\mu)\,M_{U_1}(a).
\]
\end{theorem}

\begin{proof}
For the two-row case, sum over \(c_2\in\operatorname{Fib}_{U_2}(a)\), and put
\(N_2=|A_{U_2}(c_2)|\).  The possible \(c_1\)'s are codewords agreeing with
\(U_1\) on at least \(a\) points of \(A_{U_2}(c_2)\).  If
\(N_2<2a-k\), then two such \(c_1\)'s would agree with each other on more than
\(k\) points, forcing equality.  Thus this punctured list has size at most
one.  If \(N_2\ge2a-k\), use the trivial bound \(M_{U_1}(a)\).  Summing these
two regimes gives the displayed two-row inequality.

For the \(\mu\)-row case, fix \((c_2,\ldots,c_\mu)\) and put
\[
  S=\bigcap_{i=2}^\mu A_{U_i}(c_i).
\]
The listed \(c_1\)'s must agree with \(U_1\) on at least \(a\) points of \(S\).
If \(|S|<2a-k\), the same unique-decoding argument gives at most one such
\(c_1\).  If \(|S|\ge2a-k\), use the bound \(M_{U_1}(a)\).  Counting the two
classes of \((c_2,\ldots,c_\mu)\) gives the recursion.
\end{proof}

\begin{remark}
The reduction is unconditional.  The desired L2 saving still requires an
L1-family input saying that the tail lists at agreement \(2a-k\) are small
after quotient-periodic contributions are budgeted.
\end{remark}

\subsection{Deep-point cap dependency split}

\textbf{Status: \textsc{Wrapper} / \textsc{Audit}.}
This result does not edit Paper D, but it explains why the headline MCA cap can
be routed locally without using the imported Crites--Stewart conversion.

\begin{proposition}[CS25-free route to the Paper D MCA cap]
\label{prop:deep-point-cap}
Assume the finite-field, coset, divisibility, and binomial hypotheses used in
Paper D's universal cap theorem.  Let \(C=\RS[F,D,k]\), \(C^+=\RS[F,D,k+1]\),
\(q=|F|\), and suppose the elementary quotient-locator fiber construction gives
a \(C^+\)-list of size
\[
  L\ge \binom{N}{\rho N+2}/|B|\ge q/k+1
\]
at the endpoint \(\delta_N=1-\rho-2/N\).  Then there is a support-wise MCA-bad
line for \(C\) with bad-slope density strictly larger than
\[
  \frac1{2k}\left(1-\frac nq\right).
\]
By support-wise MCA monotonicity, the same lower bound holds for all
\(\delta_N\le\delta<1-\rho\).
\end{proposition}

\begin{proof}
Let \(U:D\to F\) be the received word with a \(C^+\)-list
\(\{P_1,\ldots,P_L\}\) at agreement \(m=k+2n/N\).  For
\(\alpha\in F\setminus D\), form
\[
  f_\alpha(x)=\frac{U(x)}{x-\alpha},
  \qquad
  g_\alpha(x)=-\frac1{x-\alpha}.
\]
For every listed \(P_i\), the slope \(z=P_i(\alpha)\) gives the codeword
\[
  Q_i(X)=\frac{P_i(X)-P_i(\alpha)}{X-\alpha}\in F[X]_{<k}
\]
on the same agreement support.  Conversely every such slope comes from a
listed \(P_i\).  Since \(m>k\), the direction \(g_\alpha\) cannot be explained
by a degree-\(<k\) polynomial on the same support, so these slopes are
support-wise MCA-bad.

Averaging over \(\Omega=F\setminus D\), two distinct degree-\(<k+1\)
polynomials collide at at most \(k\) points of \(\Omega\).  Therefore some
\(\alpha\) has at least
\[
  M\ge \frac{L}{1+k(L-1)/(q-n)}
\]
distinct values among \(P_1(\alpha),\ldots,P_L(\alpha)\).  Its bad-slope
density is at least
\[
  \frac{L(q-n)}{q(q-n+k(L-1))}.
\]
This is strictly greater than \((q-n)/(2kq)\) exactly when
\[
  2kL>q-n+k(L-1).
\]
The lower bound \(L\ge q/k+1\) gives \(kL\ge q+k\), so the left-minus-right
quantity is at least \(n+2k>0\).  This proves the endpoint bound, and
monotonicity in \(\delta\) extends it to the stated interval.
\end{proof}

\subsection{Common code-line residual budget}

\textbf{Status: \textsc{New-local} finite MDS theorem.}
This M2 lemma is a guardrail for line-decoding imports: a close common
code-line exception still pays for residual coordinates outside the common
support.

\begin{proposition}[common code-line residual budget]
\label{prop:common-residual}
Let \(C\subseteq F^D\) be an MDS linear code of dimension \(k\), so every
nonzero codeword has fewer than \(k\) zeros on \(D\).  Let
\(\ell_z=f+zg\) be a received line and suppose that for some codewords
\(c_f,c_g\in C\) there is a common support \(S_0\subseteq D\) of size \(b\)
such that \(f=c_f\) and \(g=c_g\) on \(S_0\).  Fix an agreement threshold
\(a\) with \(a+b-n\ge k\).  Put \(\Omega=D\setminus S_0\),
\[
  f'=f-c_f,\qquad g'=g-c_g,\qquad
  h=\max(1,a-b),
\]
and
\[
  c_0=\#\{x\in\Omega:f'(x)=g'(x)=0\}.
\]
Every support-wise noncontained slope at agreement \(a\) satisfies
\[
  \#\{x\in\Omega:f'(x)+zg'(x)=0\}\ge h.
\]
Consequently, if \(h>c_0\), then the number of such slopes is at most
\[
  \left\lfloor\frac{|\Omega|-c_0}{h-c_0}\right\rfloor.
\]
\end{proposition}

\begin{proof}
If \(\ell_z\) is explained on a support \(S\) of size at least \(a\), then
\[
  |S\cap S_0|\ge a-|\Omega|=a+b-n\ge k.
\]
On \(S\cap S_0\), the explaining codeword agrees with \(c_f+zc_g\).  By the
MDS zero-rigidity condition, the two codewords are equal.  Thus, outside
\(S_0\), the line can only be explained where \(f'+zg'=0\), and at least
\(h=\max(1,a-b)\) such outside zeros are needed for a noncontained support.

Each non-common residual coordinate \(x\in\Omega\) with \((f'(x),g'(x))\ne
(0,0)\) can vanish for at most one finite slope \(z\).  A bad slope needs at
least \(h-c_0\) private non-common residual zeros after the \(c_0\) common zeros
are counted.  Charging these private coordinates to slopes gives the stated
packing bound.
\end{proof}

\section{Recent PR Triage: L1, F1, M1, and Challenge-Map Ledgers}
\label{sec:recent-pr-triage}

\textbf{Status: \textsc{Audit} / \textsc{Experimental} / bounded
\textsc{Proved} notes.}
The open PR batch of June 27, 2026, adds useful proof infrastructure but no new
unconditional prize-scale threshold beyond the tangent-floor consequences already
in the high-agreement ledger.  The promotion rule used here is conservative:
bounded toy theorems and exact normal forms may enter the experimental notes,
while computation-dependent numerators and generated proof programs remain
labelled as such.

\subsection{L1: proper-subgroup and monomial-prefix lanes}

Scott Hughes' \(d=3\) proper-subgroup note records the actual \(H^{2k}\)
collision object for the odd moments \(1,3,5\).  For \(p>5\), \(H\le\F_p^*\)
of order \(n\), and
\[
  S_H(a_1,a_3,a_5)
  =
  \sum_{h\in H}\psi(a_1h+a_3h^3+a_5h^5),
\]
the exact Fourier identity is
\[
  V_{H,k}^{(3)}
  =
  p^{-3}\sum_{a_1,a_3,a_5\in\F_p}
  |S_H(a_1,a_3,a_5)|^{2k}.
\]
After splitting frequencies into the linear, cubic, and quintic strata, and
using the stated one-variable Gauss/Katz inputs, the note gives
\[
0\le
V_{H,k}^{(3)}-\frac{n^{2k}}{p^3}
\le
\frac{p-1}{p^3}p^k
+\frac{p(p-1)}{p^3}(3\sqrt p)^{2k}
+\frac{p^2(p-1)}{p^3}(5\sqrt p)^{2k}.
\]
This is useful L1 infrastructure, not a reserve-scale arbitrary-word locator
local limit.

Vadim Avdeev's monomial-prefix note settles one bounded toy lane over
\[
  F=\F_{17}[z]/(z^{32}-3),\qquad H=\langle z\rangle,\qquad |H|=512,
\]
with degree convention \(\deg P\le256\).  In the range \(258\le A\le512\),
the admissible monomial agreement sizes are
\[
\begin{gathered}
258,259,260,262,264,268,272,280,\\
288,304,320,352,384,512.
\end{gathered}
\]
The proof combines a local length-\(16\) imbalance check over \(\F_{17}\), a
dyadic power-sum descent gate, and quotient-complement survivor families.  This
should be cited only as a monomial-prefix benchmark, not as an arbitrary-word
list theorem.

\subsection{F1: arbitrary-anchor residue clouds}

The F1 arbitrary-anchor update splits the balanced residue-cloud problem from
the residual-list problem.  If \(t<\sigma\), residue-line bad slopes inject into
an ordinary list for \(\RS[F,D,k+t]\) at residual slack \(\sigma-t\).  In the
balanced case \(t=\sigma\), arbitrary anchors have a real free-packing floor:
\[
  \Lambda_{\rm NC}
  \ge
  \left\lfloor\frac{|D|-k}{\sigma}\right\rfloor.
\]
Moreover this is the exact ceiling for ordered \(k\)-degenerate support
families, so any larger arbitrary-anchor packet must use dense overlaps.  The
dense-overlap gate is explicit: for two supports \(S,T\) with overlap \(J\),
compatibility of distinct slopes is controlled by whether the residue direction
\(N/E\) is degree-\(<k\) on \(J\).  This is a useful F1 warning and search
normal form, not an extension-line lift theorem.

\subsection{M1: all-line Hankel aperiodic packet}

The large M1 PR is not promoted wholesale.  The useful distillation is the exact
Hankel-pencil object.  For a line \((f,g)\), syndromes \(u=\operatorname{Syn}(f)\)
and \(v=\operatorname{Syn}(g)\), and a split complement \(T\) with locator vector
\(\ell_T\), finite bad slopes are governed by
\[
  \bigl(H_{t,j}(u)+zH_{t,j}(v)\bigr)\ell_T=0,
  \qquad
  H_{t,j}(v)\ell_T\ne0.
\]
In the \(t=2\) window this becomes
\[
  b_T\ne0,\qquad \det[a_T\ b_T]=0,
  \qquad
  a_T=H_{2,j}(u)\ell_T,\quad b_T=H_{2,j}(v)\ell_T.
\]
Same-slope one-exchange collisions lie in fixed-slope root slices, and the
residual one-exchange graph can be further organized into quadratic companion
slices and top packets.  These reductions define a promising M1 proof program,
but they do not yet prove the desired uniform polynomial bound on residual
aperiodic slope images.

\subsection{M1: same-slope one-exchange root slices}

The first local reduction in the all-line Hankel packet is elementary and
proved in \path{experimental/notes/m1/m1_same_slope_root_slice_lemma.md}.  Let
\[
  L_z=H_{t,j}(u)+zH_{t,j}(v)
\]
be the finite-slope landing map.  Fix a \((j-1)\)-core \(R\), and write
\[
  T_y=R\cup\{y\},\qquad \ell_{T_y}=(X-y)\ell_R .
\]
If two distinct one-exchange extensions \(T_{y_1}\) and \(T_{y_2}\) satisfy
\[
  L_z\ell_{T_{y_1}}=L_z\ell_{T_{y_2}}=0,
\]
then subtracting gives
\[
  (y_2-y_1)L_z\ell_R=0.
\]
Thus \(L_z\ell_R=0\), and substituting into either endpoint equation gives
\[
  L_z(X\ell_R)=0.
\]
Consequently
\[
  L_z\ell_{T_y}=0\qquad\text{for every }y.
\]
So every same-slope one-exchange edge lies in the full fixed-slope root slice
with core \(R\).  After fixed-slope root slices have been charged, the residual
one-exchange graph has only different-slope edges.  This is a local reduction,
not a bound for the remaining different-slope codegree.  The verifier
\path{experimental/scripts/verify_m1_same_slope_root_slice_lemma.py} checks the
subtraction identity and its row-wise linear-map consequence in small fields.

The same two padded equations also lift the root slice to a higher-slack
Hankel core.  With \(w_z=u+zv\), the equations
\[
  H_{t,j}(w_z)(\ell_R,0)^T=0,\qquad
  H_{t,j}(w_z)(0,\ell_R)^T=0
\]
are exactly the first \(t\) and last \(t\) rows of
\[
  H_{t+1,j-1}(w_z)\ell_R=0.
\]
Thus a same-slope root slice is not a fresh \(t\)-level residual multiplicity:
it belongs to the lifted \((t+1,j-1)\) Hankel-core ledger at the same finite
slope.

The same mechanism has a two-exchange full-plane form.  For a \((j-2)\)-core
\(R\), write a two-root locator in elementary coordinates
\[
  \ell_{T_{s,p}}=(X^2-sX+p)\ell_R,\qquad s=x+y,\quad p=xy .
\]
For fixed finite slope \(z\),
\[
  L_z\ell_{T_{s,p}}
  =
  L_z(X^2\ell_R)-sL_z(X\ell_R)+pL_z(\ell_R)
\]
is affine-linear in \((s,p)\).  Hence three affinely non-collinear same-slope
points in the two-root plane force
\[
  L_z\ell_R=L_z(X\ell_R)=L_z(X^2\ell_R)=0,
\]
equivalently
\[
  H_{t+2,j-2}(u+zv)\ell_R=0.
\]
Thus full two-exchange planes are charged to the lifted
\((t+2,j-2)\) Hankel ledger.  After fixed-root lines and full planes are
removed, any residual same-slope two-exchange component through \(R\) is a
line packet in the \((s,p)\)-plane.

The full-plane statement is the \(h=2\) case of a general elementary-packet
lift.  For a \((j-h)\)-core \(R\), write
\[
  \ell_{R,c}=(X^h+c_{h-1}X^{h-1}+\cdots+c_0)\ell_R,
  \qquad c\in F^h.
\]
For fixed finite slope \(z\), the map \(c\mapsto L_z\ell_{R,c}\) is
affine-linear.  Hence if \(h+1\) affinely independent coefficient points are
killed at the same slope, then
\[
  L_z(X^m\ell_R)=0,\qquad 0\le m\le h.
\]
These \(h+1\) padded equations are exactly the row blocks of
\[
  H_{t+h,j-h}(u+zv)\ell_R=0.
\]
Thus every full affine-rank same-slope \(h\)-exchange packet is charged
losslessly to the lifted \((t+h,j-h)\) Hankel-core ledger.  Applying the same
argument separately to \(u\) and \(v\) gives the simultaneous-kernel version:
full affine-rank packets in \(K_{r,d}(u,v)\) lift to \(K_{r+h,d-h}(u,v)\).
The lower-rank case is also exact.  If \(C\subset F^h\) is any killed
coefficient packet and \(\operatorname{aff}(C)=c_*+W\), then
\[
  L_z(X^h\ell_R)+\sum_m c_{*,m}L_z(X^m\ell_R)=0,\qquad
  \sum_m w_mL_z(X^m\ell_R)=0\quad(w\in W).
\]
Conversely these equations kill the whole affine subpacket \(c_*+W\).  Thus
after full-rank packets are charged, the residual same-slope elementary
packets are precisely lower-rank affine subpackets; for \(h=2\) these are the
lines classified next.
Among coefficient hyperplanes, the fixed-root ones are exactly the evaluation
hyperplanes.  Namely
\[
  b+\sum_{m=0}^{h-1}a_m c_m=0
\]
is the locus \(P_c(\alpha)=0\) for a fixed finite root \(\alpha\) iff, up to a
nonzero scalar, \((a_0,\ldots,a_{h-1},b)=(1,\alpha,\ldots,\alpha^h)\).  Such
hyperplanes are charged to the \((h-1)\)-exchange root-slice ledger.  For
\(h=2\), this is \(p-\alpha s+\alpha^2=0\), equivalently
\((x-\alpha)(y-\alpha)=0\).
More generally, every coefficient hyperplane has a one-root fiber dichotomy.
Fix a monic \((h-1)\)-root factor
\[
  Q_d(X)=X^{h-1}+d_{h-2}X^{h-2}+\cdots+d_0
\]
and write \(d_{h-1}=1,d_{-1}=0\).  The extension
\[
  P_y=(X-y)Q_d
\]
has coefficients \(c_m(y)=d_{m-1}-yd_m\), so the hyperplane equation restricts
to
\[
  b+\sum_m a_m c_m(y)
  =
  \Bigl(b+\sum_m a_m d_{m-1}\Bigr)-y\sum_m a_m d_m .
\]
Thus a fixed \((h-1)\)-core has at most one extension in the hyperplane unless
the whole affine one-root line lies there.  In the killed same-slope packet
setting, such a full line is exactly a one-exchange root slice and lifts to the
\((t+1,j-1)\) Hankel core.  After these full one-root fibers are charged, a
codimension-one coefficient-hyperplane packet has no residual one-exchange
edges.
The same 0/1/all dichotomy holds for any affine rank-defect packet
\(A=c_*+W\subset F^h\).  For the same one-root fiber, put
\[
  d_{\rm vec}=(d_0,\ldots,d_{h-1}).
\]
If two distinct roots \(y_1,y_2\) have \(c(y_i)\in A\), then
\[
  c(y_2)-c(y_1)=-(y_2-y_1)d_{\rm vec}\in W,
\]
so \(d_{\rm vec}\in W\) and the whole affine line \(c(y)\) lies in \(A\).
Thus, after full one-root fibers are root-slice charged, no residual affine
rank-defect packet has one-exchange edges.
The two-root fiber has the same affine form.  For a fixed monic
\((h-2)\)-root factor \(Q_e\), write
\[
  P_{s,p}=(X^2-sX+p)Q_e,\qquad
  c_m(s,p)=e_{m-2}-s e_{m-1}+p e_m .
\]
If an affine rank-defect packet contains three non-collinear points of this
\((s,p)\)-plane, then both plane directions lie in \(W\), so the whole
two-root plane is contained in the packet and is charged to the
\((t+2,j-2)\) full-plane lift.  After full planes are charged, every residual
two-root fiber intersection is at most a line packet.
In general, for a moving monic \(r\)-root factor \(B_a\) multiplying a fixed
\((h-r)\)-root factor \(Q\), the map \(a\mapsto B_aQ\) is an affine embedding
of \(F^r\) into coefficient space.  The preimage of \(A=c_*+W\) is therefore
an affine subspace of \(F^r\): if it has full affine rank \(r\), the whole
moving \(r\)-root fiber is contained in \(A\) and is charged to the
\((t+r,j-r)\) full-packet lift; otherwise the residual intersection has
dimension at most \(r-1\).
Over a finite field of size \(Q_F\), this gives the formal residual count
\[
  \#\{a\in F^r:\operatorname{coeff}(B_aQ)\in A\}\le Q_F^{r-1}
\]
on every fixed \(r\)-root fiber after the full-fiber charge; split, domain,
quotient, tangent, and noncontainment filters only shrink this count.
Consequently the residual \(r\)-exchange graph inside one affine packet has
maximum degree at most
\[
  \binom{h}{r}\bigl(Q_F^{r-1}-1\bigr),
\]
because each vertex has only \(\binom hr\) fixed \((h-r)\)-subfactors to
complete.

The residual two-root lines are exactly the expected
models: if
\[
  As+Bp+C=0
\]
and \(B=0\), then \(x+y=s_0\) is a fixed-sum packet; if \(B\ne0\), then after
writing \(c=-A/B\), \(\beta=-C/B\), the equation becomes
\[
  (x-c)(y-c)=c^2+\beta .
\]
The zero right-hand side is a fixed-root line, and the nonzero right-hand side
is a product-Mobius packet.  Hence after fixed-root lines are charged, the
only residual line packets are fixed-sum and nondegenerate product-Mobius
packets.

In the actual \(t=2\) Hankel window, these non-fixed line packets cannot be
constant-slope residuals.  If
\[
  H_{2,j}(u+zv)\ell_{R,s,p}
  =(d_2-sd_1+pd_0,\ d_3-sd_2+pd_1)
\]
vanishes identically on a fixed-sum line \(s=s_0\), then the coefficient of
\(p\) gives \((d_0,d_1)=0\), and the Hankel overlap forces
\(d_2=d_3=0\).  If it vanishes identically on a nondegenerate
product-Mobius line \(p=cs-c^2+\mu\), \(\mu\ne0\), then
\((d_1,d_2)=c(d_0,d_1)\) and
\((d_2,d_3)=(c^2-\mu)(d_0,d_1)\), so \(\mu d_0=0\) and again all four
\(d_i\) vanish.  Thus every constant-slope non-fixed line packet is already
charged to the full-plane lift \(H_{4,j-2}(u+zv)\ell_R=0\).  After fixed-root
and full-plane charges, the active finite-slope map on each surviving
fixed-sum or product-Mobius packet is injective.

In the \(t=2\) Hankel window the same one-exchange core has a quadratic
determinant gate.  Put
\[
  a_y=H(u)\ell_{T_y}=a_X-y a_0,\qquad
  b_y=H(v)\ell_{T_y}=b_X-y b_0
  \qquad\text{in }F^2.
\]
Then the finite-slope determinant condition is
\[
  \Delta_R(y)=\det(a_y,b_y)=0,\qquad b_y\ne0,
\]
where
\[
  \Delta_R(y)
  =
  \det(a_X,b_X)
  -y\bigl(\det(a_0,b_X)+\det(a_X,b_0)\bigr)
  +y^2\det(a_0,b_0).
\]
Thus a nonzero \(\Delta_R\) admits at most two anchors \(y\).  If three
distinct anchors pass the determinant gate, then \(\Delta_R\equiv0\), so the
core is a ruled determinant branch.  The ruled branch first has the abstract
affine dichotomy: if
\[
  \det(a_X-y a_0,\ b_X-y b_0)\equiv0,
\]
then either there is a fixed finite slope \(z_0\) with
\[
  a_y+z_0 b_y=0\qquad\text{for every }y,
\]
or \(b_y=0\) for every \(y\), or all four vectors
\[
  a_X,\ a_0,\ b_X,\ b_0
\]
lie in a single output line.  In the last case write
\[
  a_y=\alpha(y)c,\qquad b_y=\beta(y)c .
\]
For arbitrary affine pencils this last case could be a moving-slope branch, but
the Hankel shift structure rules it out here.  If
\[
  c_i(w)={\rm row}_i(H_{3,j-1}(w)\ell_R),\qquad i=0,1,2,
\]
then
\[
  H_{2,j}(w)\ell_{T_y}
  =
  (c_1(w)-yc_0(w),\ c_2(w)-yc_1(w)).
\]
A common output line imposes two recurrence equations on
\((c_0(w),c_1(w),c_2(w))\), whose solution space is one-dimensional.  Applying
this to \(w=u\) and \(w=v\) makes the two Hankel pencils proportional, unless
\(b_y=0\) identically.  Therefore every active ruled Hankel core is already a
fixed finite-slope core; after fixed-slope root slices are charged, ruled cores
contribute no residual one-exchange edges.

The same note also records the exact one-exchange triangle dichotomy in the
Johnson graph.  Three \(j\)-set locators that are pairwise one-exchange are
either a star
\[
  T_i=R\cup\{y_i\},\qquad |R|=j-1,
\]
or a top-packet triangle
\[
  T_i=U\setminus\{x_i\},\qquad |U|=j+1.
\]
Thus every active star triangle is precisely a three-anchor event through one
\((j-1)\)-core, hence a ruled determinant core by the quadratic gate above.
Since ruled Hankel cores collapse to fixed-slope or inactive cores, residual
one-exchange triangles after fixed-slope root-slice charging are therefore
top-packet triangles.

Top packets then lift losslessly to a \(t=1\) Hankel kernel.  If
\[
  T_x=U\setminus\{x\},\qquad \ell_U=(X-x)\ell_{T_x},
\]
then for every syndrome vector \(w\),
\[
  H_{1,j+1}(w)\ell_U
  =
  {\rm row}_1\bigl(H_{2,j}(w)\ell_{T_x}\bigr)
  -x\,{\rm row}_0\bigl(H_{2,j}(w)\ell_{T_x}\bigr).
\]
Thus a \(t=2\) incidence for \(T_x\) at slope \(z\) implies
\[
  (H_{1,j+1}(u)+zH_{1,j+1}(v))\ell_U=0.
\]
If two members of the same top packet have distinct slopes, subtracting the
two lifted scalar equations gives
\[
  H_{1,j+1}(u)\ell_U=H_{1,j+1}(v)\ell_U=0.
\]
Therefore every residual top-packet edge after fixed-slope charging lies over a
common lifted \(t=1\) kernel.  This turns the top-packet triangle residual into
a higher-slack kernel target rather than a fresh one-exchange multiplicity
phenomenon.

Writing
\[
  K_{\rm top}(u,v)
  =
  \{\,U\subset D:\ |U|=j+1,\
      H_{1,j+1}(u)\ell_U=H_{1,j+1}(v)\ell_U=0\,\},
\]
the top-packet branch has the explicit compression ledger
\[
  E_{\rm top}(A_{\rm res})
  \le \binom{j+1}{2}|K_{\rm top}(u,v)|,\qquad
  {\rm Tri}_1(A_{\rm res})
  \le \binom{j+1}{3}|K_{\rm top}(u,v)|.
\]
Indeed a residual top edge recovers its union \(U\), and (because residual
one-exchange edges have distinct slopes) the lift above forces
\(U\in K_{\rm top}(u,v)\).  Inside such a packet,
\[
  H_{2,j}(w)\ell_{U\setminus\{x\}}
  =
  \rho_x(w)(1,x),\qquad w\in\{u,v\},
\]
so the remaining slope label is the scalar ratio
\(\rho_x(u)+z\rho_x(v)=0\), \(\rho_x(v)\ne0\).  Thus top-packet edges and
triangles are charged to the lifted \(t=1\) top-kernel family.

The top-kernel family also has an exact root-slice recursion.  More generally,
for \(r\ge1\) and locator size \(d\), put
\[
  K_{r,d}(u,v)
  =
  \{\,U\subset D:\ |U|=d,\
      H_{r,d}(u)\ell_U=H_{r,d}(v)\ell_U=0\,\}.
\]
If \(U_{y_i}=R\cup\{y_i\}\in K_{r,d}(u,v)\) for two distinct anchors through
the same \((d-1)\)-core \(R\), then the same subtraction argument applied
separately to \(u\) and \(v\) gives
\[
  H_{r,d}(w)\ell_R=H_{r,d}(w)(X\ell_R)=0,
  \qquad w\in\{u,v\}.
\]
These are exactly the first \(r\) and last \(r\) rows of
\[
  H_{r+1,d-1}(w)\ell_R=0.
\]
Hence \(R\in K_{r+1,d-1}(u,v)\).  So one-exchange collisions inside the lifted
top-kernel ledger are charged losslessly to the next simultaneous Hankel
kernel; after those root slices are paid, the residual \(K_{r,d}\) family has
no one-exchange edges.  This records the residual-depth frontier shift as an
identity rather than a multiplicative loss source.

The same Hankel shift also classifies the one-outside external-anchor ruled
branch from the boundary-off audit.  For a shadow \(S\subset D\), \(|S|=j-1\),
and an external anchor \(\beta\notin D\),
\[
  \ell_{S,\beta}=(X-\beta)\ell_S .
\]
The \(t=2\) determinant gate is again quadratic in \(\beta\).  If it is
nonzero, at most two external anchors pass.  If it vanishes identically, then
with
\[
  c_i(w)={\rm row}_i(H_{3,j-1}(w)\ell_S)
\]
one has
\[
  H_{2,j}(w)\ell_{S,\beta}
  =
  (c_1(w)-\beta c_0(w),\ c_2(w)-\beta c_1(w)).
\]
The same recurrence argument as above forces the ruled external-anchor branch
to be inactive or fixed finite slope.  In the fixed-slope case it is a boundary
root slice and lifts to
\[
  H_{3,j-1}(u+z_0v)\ell_S=0.
\]
Thus, after fixed-slope boundary slices are charged, each shadow \(S\) supports
at most two active non-ruled external anchors.  This classifies the ruled
\({\rm Boundary}_{\rm off}\) branch locally, but it is not yet a global bound
for the one-outside target image.
Equivalently, projecting a residual one-outside target
\(B_\beta=S\cup\{\beta\}\) to its domain shadow \(S\) has fibers of size at
most two.  If two external anchors over the same shadow had the same finite
slope \(z\), then subtraction gives the lifted boundary root slice
\[
  H_{3,j-1}(u+zv)\ell_S=0,
\]
so after that charge the remaining two-anchor shadow fibers have distinct
finite slopes.  Hence
\[
  |{\rm Boundary}_{\rm off}^{\rm res}|
  \le 2|{\rm Shadow}_{\rm off}^{\rm res}|.
\]
This reduces the live one-outside object to a shadow-image problem rather than
an external-anchor multiplicity problem.
The shadow condition itself is intrinsic.  Writing
\[
  c_i={\rm row}_i(H_{3,j-1}(u+zv)\ell_S),\qquad 0\le i\le2,
\]
there is an external anchor \(\beta\) with
\(H_{2,j}(u+zv)\ell_{S,\beta}=0\) iff either
\((c_0,c_1,c_2)=(0,0,0)\), the already charged lifted boundary slice, or
\[
  c_0\ne0,\qquad c_1^2=c_0c_2,\qquad \beta=c_1/c_0 .
\]
Thus the residual boundary shadow image lies in a rank-one scalar Hankel locus
with recovered anchor \(c_1/c_0\notin D\), plus the usual active filter.
As \(z\) varies, write \(c_i(z)=a_i+zb_i\).  The recovered-anchor gate is the
quadratic
\[
  Q_S(z)=c_1(z)^2-c_0(z)c_2(z).
\]
If \(Q_S\) is nonzero, the fixed shadow contributes at most two candidate
finite slopes and hence at most two recovered anchors.  If \(Q_S\equiv0\), the
line \((c_0,c_1,c_2)(z)=a+zb\) lies on one generator of the rank-one cone
\(x_1^2=x_0x_2\).  Thus any finite recovered anchor is independent of \(z\),
and both \(H_{2,j}(u)\ell_{S,\beta}\) and
\(H_{2,j}(v)\ell_{S,\beta}\) vanish there.  The active filter removes this
degenerate branch, while zero triples are already charged to the lifted
boundary slice.  The remaining boundary-shadow problem is therefore a nonzero
quadratic-root problem over the shadows.
Equivalently, with \(r_\beta=(1,\beta,\beta^2)\) and
\(h_\beta(c)=(c_1-\beta c_0,c_2-\beta c_1)\), define the anchor gate
\[
  A_S(\beta)=\det(a,b,r_\beta).
\]
Then \(A_S(\beta)=\det(h_\beta(a),h_\beta(b))\).  If
\(h_\beta(b)\ne0\), a root \(A_S(\beta)=0\) gives a unique finite slope
\(z\) with \(h_\beta(a+zb)=0\); nonzero triples are exactly the recovered
rank-one anchors above, while zero triples are charged.  If \(A_S\equiv0\),
then \(a,b\) are proportional and the branch is inactive or charged.  Thus the
same residual target is a nondegenerate conic-secant incidence: a nonzero
quadratic in the outside anchor \(\beta\), with the active filter imposed.
Fixing \(\beta\) gives a further one-root reduction.  For
\(|R|=j-2\), put \(P_{R,\beta}=(X-\beta)\ell_R\) and
\(\ell_{R,\beta,y}=(X-y)P_{R,\beta}\).  The determinant
\[
  \det(H_{2,j}(u)\ell_{R,\beta,y},
       H_{2,j}(v)\ell_{R,\beta,y})
\]
is quadratic in \(y\).  A nonzero determinant has at most two domain roots;
an identically zero determinant is inactive or fixed finite slope, and the
fixed-slope case lifts to
\[
  H_{3,j-1}(u+zv)P_{R,\beta}=0 .
\]
After charging these one-outside lifted boundary-core slices, projection
\((R\cup\{y\},\beta)\mapsto(R,\beta)\) has residual fibers of size at most two.
Combining this with the shadow-fiber bound, for each fixed \((j-2)\)-core
\(R\) the residual boundary incidences between external anchors \(\beta\) and
domain roots \(y\in D\setminus R\) form a bipartite graph of degree at most two
on both sides.  Hence each fixed \(R\) supports at most \(2(n-j+2)\) residual
external anchors.  Thus the next boundary object is the lower-dimensional
boundary-core image, or still further the active \((j-2)\)-core image.
For fixed \(R\), write
\[
  U_m=H_{2,j}(u)(X^m\ell_R),\qquad
  V_m=H_{2,j}(v)(X^m\ell_R),\qquad 0\le m\le2 .
\]
Then the remaining fixed-core incidence is cut out by the symmetric bidegree
\((2,2)\) determinant
\[
  \Delta_R(\beta,y)=
  \det\!\left(
    U_2-(\beta+y)U_1+\beta y U_0,\,
    V_2-(\beta+y)V_1+\beta y V_0
  \right).
\]
The fixed-shadow and fixed-anchor quadratics are the two coordinate
specializations of this same determinant.  Thus the active core image is a
concrete bidegree-two incidence target, not an uncontrolled external-anchor
search.
Equivalently, \(\Delta_R(\beta,y)=F_R(\beta+y,\beta y)\), where
\[
  F_R(s,p)=\det(U_2-sU_1+pU_0,\;V_2-sV_1+pV_0)
\]
is the ordinary two-root elementary determinant.  Fixing either coordinate is
the fixed-root line \(p=\alpha s-\alpha^2\) in the \((s,p)\)-plane, so the
boundary-core target is a mixed-domain slice of the same determinant geometry
used by the fixed-root, fixed-sum, and product-Mobius line-packet ledger.
Moreover the mixed slice of each surviving non-fixed line packet is explicit:
on a fixed-sum line \(x+y=s_0\) one has \(\beta=s_0-y\), while on a
nondegenerate product-Mobius line \((x-c)(y-c)=\mu\ne0\) one has
\(\beta=c+\mu/(y-c)\).  Thus, over a fixed core \(R\), such a packet
contributes at most \(n-j+2\) one-outside boundary-core pairs, namely the
domain roots whose packet partner leaves \(D\); the \(\mu=0\) case is the
fixed-root line already charged.
The same elementary-plane form also controls slope multiplicity on the
remaining conic branch.  For fixed finite slope \(z\), put
\(W_m=U_m+zV_m\).  The same-slope equations are
\[
  W_2-sW_1+pW_0=0,
\]
a pair of affine-linear equations in \((s,p)\).  Hence the same-slope fiber is
empty, a point, an affine line, or the full plane.  A repeated slope on two
distinct elementary points therefore forces a constant-slope affine line, or
the full-plane lift.  Fixed-root lines are already boundary/root-slice charges,
and non-fixed constant-slope lines collapse to the full-plane charge above.
Thus after these charges the finite-slope label is injective on the live
fixed-core boundary conic, so the quartic point count below also controls the
number of live slopes.
For the complementary conic case, write
\[
F_R(s,p)=f_{20}s^2+f_{11}sp+f_{02}p^2+f_{10}s+f_{01}p+f_{00}.
\]
Then for fixed \(y\),
\[
F_R(\beta+y,\beta y)=A_R(y)\beta^2+B_R(y)\beta+C_R(y),
\]
where
\[
\begin{aligned}
A_R(y)&=f_{20}+f_{11}y+f_{02}y^2,\\
B_R(y)&=f_{10}+(2f_{20}+f_{01})y+f_{11}y^2,\\
C_R(y)&=f_{00}+f_{10}y+f_{20}y^2 .
\end{aligned}
\]
In odd characteristic the nondegenerate anchor count is controlled by the
quartic discriminant \(B_R(y)^2-4A_R(y)C_R(y)\); a fully zero quadratic is the
fixed-root component \(p=y s-y^2\), already charged.  Thus the remaining
non-line boundary-core target is a quartic discriminant/Kummer trace over
domain roots, before outside-domain and active filters.  If this quartic
discriminant is identically zero, then coefficient comparison shows that
either \(F_R=\lambda L^2\) for an affine line \(L\), already a fixed-sum,
fixed-root, or product-Mobius line-packet branch, or
\(F_R=\lambda(s^2-4p)\).  The latter restricts to
\(\lambda(\beta-y)^2\) on the boundary slice, hence contributes only
\(\beta=y\), which is removed by the outside-domain/distinct-root filters.
So the live conic target after line-packet and fixed-root charges has nonzero
quartic discriminant.  If \(A_R\not\equiv0\), then away from at most two roots
of \(A_R\) the quadratic formula identifies anchor roots with the double cover
\[
  W_R^2=B_R(y)^2-4A_R(y)C_R(y),\qquad
  W_R=2A_R(y)\beta+B_R(y).
\]
After squarefree reduction this cover has genus at most one, since the
discriminant has degree at most four.  If \(A_R\equiv0\), the determinant is
affine-linear in \((s,p)\) and has already returned to the line-packet branch.
Thus the live non-line boundary target has uniformly bounded genus/conductor,
up to at most two linear exceptional fibers.
There is also an exact full-subgroup Kummer gate.  If \(D\le\F_p^*\) has index
\(e\), \(\chi\) has kernel \(D\), and \(\ell=\operatorname{lcm}(e,2)\), then
away from the charged fixed-root fibers and the at most two \(A_R=0\)
exceptions the full-subgroup count expands as
\[
  N_R(D)
  = |D|+\frac1e\sum_{a=0}^{e-1}
      \sum_{y\in\F_p^*}\chi^a(y)\chi_2(\operatorname{Disc}_R(y))+O(1).
\]
Combining \(\chi\) and \(\chi_2\) through a character of order \(\ell\), the
\(a\)-term is geometrically trivial only if
\[
  \operatorname{div}\!\left(y^{a\ell/e}\operatorname{Disc}_R(y)^{\ell/2}\right)
  \equiv 0 \pmod{\ell}.
\]
Equivalently, all nonzero roots of \(\operatorname{Disc}_R\) have even multiplicity and the
zero-root parity matches \(a\).  Hence the squarefree part of a
no-cancellation quartic is forced to be either \(1\) or \(y\):
\[
  \operatorname{Disc}_R(y)=cG(y)^2
  \qquad\hbox{or}\qquad
  \operatorname{Disc}_R(y)=c\,y\,G(y)^2 .
\]
The first case forces \(a=0\), while the second forces \(e\) even and
\(a=e/2\).  If this power gate fails, the standard one-dimensional
Kummer-Weil bound on \(\mathbb P^1\) has support contained in the four
discriminant roots plus \(0\) and \(\infty\), so the term has
depth-independent conductor and is bounded by \(4\sqrt p\).  Thus any
no-cancellation boundary quartic is an explicit rational graph or fixed
square-root rational graph branch, not a hidden growing-conductor family.
It follows that the full-subgroup live count over a fixed core satisfies
\[
  N_R^{\rm live}(D)\le |D|+4\sqrt p+O(1)
\]
when the power gate is absent, and
\[
  N_R^{\rm live}(D)\le 2|D|+4\sqrt p+O(1)
\]
in general.  There is at most one geometrically trivial Fourier term: the
principal term in the \(cG^2\) case or the quadratic subgroup term in the
\(cyG^2\) case.  The outside-domain and active filters only delete points,
and the same-slope fiber injection above turns this point count into the same
live slope count for the fixed core.
Combining with the degree-two graph bound gives the ledger-facing non-line
conic estimate
\[
  N_R^{\rm conic}(D)
  \le \min\{2(n-j+2),\, |D|+4\sqrt p+O(1)\}
\]
when the power gate is absent, and
\[
  N_R^{\rm conic}(D)
  \le \min\{2(n-j+2),\, 2|D|+4\sqrt p+O(1)\}
\]
in general.  Reducible fixed-sum and product-Mobius line packets remain in
the separate graph/packet ledger; the non-line fixed-core geometry has no
remaining uncontrolled multiplicity.
The same fixed-core target also has a slope-side recurrence chart.  Write
\[
  a_i={\rm row}_i(H_{4,j-2}(u)\ell_R),\qquad
  b_i={\rm row}_i(H_{4,j-2}(v)\ell_R),\qquad
  c_i(z)=a_i+zb_i .
\]
For \(P(X)=X^2-sX+p\), the equation
\[
  H_{2,j}(u+zv)(P\ell_R)=0
\]
is exactly
\[
  c_2-sc_1+pc_0=0,\qquad c_3-sc_2+pc_1=0 .
\]
If \(Q_R(z)=c_0c_2-c_1^2\ne0\), then
\[
  s=\frac{c_0c_3-c_1c_2}{Q_R(z)},\qquad
  p=\frac{c_1c_3-c_2^2}{Q_R(z)},
\]
and the split-root condition is governed by the quartic numerator
\[
  (c_0c_3-c_1c_2)^2
  -4(c_0c_2-c_1^2)(c_1c_3-c_2^2).
\]
If \(Q_R(z)=0\) and a solution exists, then the four \(c_i\)'s are either all
zero, giving the full-plane lift, or form a geometric progression
\(c_i=c_0\alpha^i\), in which case every solution has \(P(\alpha)=0\) and is
a fixed-root line.  Thus the residual non-line root-core target lives in the
nonzero-denominator recurrence chart.
Set
\[
  A_R=c_0c_3-c_1c_2,\qquad B_R=c_1c_3-c_2^2,\qquad
  \Theta_R=A_R^2-4Q_RB_R .
\]
Because the \(c_i\)'s are linear in \(z\), the polynomials
\(Q_R,A_R,B_R\) have degree at most two and \(\Theta_R\) has degree at most
four.  On \(Q_R\ne0\), a split lift with roots \(r_1,r_2\) gives
\[
  y=Q_R(z)(r_1-r_2),\qquad y^2=\Theta_R(z),
\]
and conversely a point on \(Y^2=\Theta_R(z)\) recovers the formal roots
\((s_R(z)\pm y/Q_R(z))/2\).  Therefore, after the denominator-zero
fixed-root/full-plane cases are charged, the residual split root-core slopes
are filtered points on a degree-four double cover of the slope line.  Its
squarefree reduction has genus at most one; if \(\Theta_R\) is a square the
roots are rational functions of \(z\), and otherwise the standard Kummer-Weil
bound has depth-independent conductor.
The domain/outside filter does not change this conductor scale.  If
\(D=(\mathbb F_p^*)^e\) and \(\chi\) is the quotient character of order \(e\),
then on the cover, away from \(Q_R=0\),
\[
  r_\pm(z,Y)=\frac{A_R(z)\pm Y}{2Q_R(z)},\qquad
  1_D(x)=\frac1e\sum_{a=0}^{e-1}\chi^a(x),
\]
with characters extended by zero at \(x=0\).  Hence the condition
\(r_+\in D\), \(r_-\notin D\) is an explicit finite sum of Kummer traces
\[
  \sum_{(z,Y)} \chi^a(r_+(z,Y))\chi^b(r_-(z,Y))
\]
on the same genus-zero/genus-one cover.  The finite poles lie over the at most
two roots of \(Q_R\); finite zeros lie over the at most two roots of \(B_R\);
the remaining support is at infinity and at the at-most-four branch points of
\(\Theta_R\).  Thus every non-power summand has \(O_e(\sqrt p)\) size by the
standard curve Kummer-Weil bound with conductor independent of depth.  The only
no-cancellation cases are explicit power-divisor congruences for
\(r_+^a r_-^b\) on this bounded divisor support.
If \(B_R\equiv0\), then \(p_R=0\) on \(Q_R\ne0\), so every formal two-root
extension has the fixed root \(0\); this is the fixed-root line \(p=0\) and is
already charged.  Otherwise the zero-root points lie over the at most two
roots of \(B_R\).  On the open cover \(C_R^\times\), where
\(Q_Rr_+r_-\ne0\), put
\[
  N_R(z)=B_R(z)/Q_R(z)=r_+(z,Y)r_-(z,Y).
\]
Then
\[
  1_D(r_+)(1-1_D(r_-))=1_D(r_+)(1-1_D(N_R(z))).
\]
Indeed, once \(r_+\in D\), the condition \(r_-\in D\) is equivalent to
\(N_R(z)\in D\).  Thus the outside-root test is a slope-line norm filter; the
remaining cover-level condition is only \(r_+\in D\).  Equivalently,
\[
  N_R^{+,-}(D)
  =\frac1e\sum_a\sum_{C_R^\times}\chi^a(r_+)(1-1_D(N_R(z)))+O(1),
\]
which is the same as the expansion below after reindexing
\(S_{a+b,b}=\sum\chi^a(r_+)\chi^b(N_R)\).
For a fixed slope with \(Q_RB_R\ne0\), the mixed ordered points inject into
the norm-outside slope set: if \(N_R(z)\in D\), then no mixed point occurs,
while if \(N_R(z)\notin D\), at most one of the two roots can lie in \(D\).
Thus
\[
  N_R^{+,-}(D)
  \le Z_R^{\rm norm-out}
  :=\sum_{Q_RB_R\ne0}\left(1-1_D(B_R/Q_R)\right).
\]
The right side has the \(\mathbb P^1\) expansion
\[
  Z_R^{\rm norm-out}
  =(1-1/e)\#\{Q_RB_R\ne0\}
  -\frac1e\sum_{a=1}^{e-1}\sum_{Q_RB_R\ne0}\chi^a(B_R/Q_R).
\]
Hence, unless \(B_R/Q_R\) is a quotient-character power, the standard
\(\mathbb P^1\) Kummer bound gives
\[
  N_R^{+,-}(D)\le (e-1)|D|+C_e\sqrt p+O_e(1).
\]
The norm-power exception is explicit: for a nonprincipal \(\chi^a\) of order
\(m=e/\gcd(e,a)\), geometric triviality is exactly
\(\mathrm{div}(B_R/Q_R)\equiv0\pmod m\).  Since \(B_R,Q_R\) have degree at
most two, \(m\ge3\) forces \(B_R/Q_R\) to be constant, while \(m=2\) says that
the reduced rational function \(B_R/Q_R\) is a square up to scalar.  Thus odd
index has only the constant-norm exception, and even index has only the
additional quadratic square-norm branch.
The constant-norm branch is already charged: \(B_R/Q_R=\gamma\) means
\(p_R(z)=\gamma\), hence every live two-root extension lies on the elementary
line \(p=\gamma\), i.e. the fixed-root line if \(\gamma=0\) and otherwise the
center-zero product-Mobius packet \(xy=\gamma\).  In the remaining nonconstant
square-norm case, say \(B_R/Q_R=\gamma M(z)^2\) and \(e\) is even, all
nonprincipal terms in the norm-outside expansion are bounded-support
\(\mathbb P^1\) Kummer sums except the quadratic coefficient
\(\chi^{e/2}\), where
\(\chi^{e/2}(B_R/Q_R)=\chi^{e/2}(\gamma)=\epsilon\).  Therefore
\[
  Z_R^{\rm norm-out}
  =\frac{e-1-\epsilon}{e}\#\{Q_RB_R\ne0\}
   +O_e(\sqrt p)+O_e(1).
\]
Thus the positive-parity square-norm branch has the genus-free bound
\((e-2)|D|+C_e\sqrt p+O_e(1)\), while the negative-parity square-norm branch
has the full slope-line main term and remains the square-norm obstruction for
the cover-level sums.
In the positive-parity branch, writing \(N_R=\gamma M^2\), the condition
\(N_R\in D\) is \(\chi(M)^2=\chi(\gamma)^{-1}\).  This cuts out two quotient
classes of the degree-one slope parameter \(M\), so the norm-outside slopes
are exactly the remaining \(e-2\) classes, up to zeros and poles.
In the negative-parity branch \(N_R\notin D\) on the open slope line, so the
mixed count collapses to
\[
  N_R^{+,-}(D)=e^{-1}|C_R^\times|+e^{-1}\sum_{a=1}^{e-1}S_a^++O(1).
\]
All nonquadratic one-root terms are then nontrivial; the only possible large
term is the explicit branch where \(r_+\) is a square.  Since \(\deg r_+\le2\),
this branch is fixed-root or rational.  Thus the negative-parity square-norm
genus-one branch has bound \(|D|+C_e\sqrt p+O_e(1)\), while the rational
one-root-square branch has a final sign: if \(r_+=\alpha H^2\), then
\(\chi^{e/2}(\alpha)=-1\) cancels the main term and
\(\chi^{e/2}(\alpha)=+1\) gives the safe bound
\(2|D|+C_e\sqrt p+O_e(1)\).  In the positive-sign case the degree-one
parameter \(h=H\) satisfies \(r_+(h)=\alpha h^2\), so
\(r_+(h)\in D\) is exactly the two-coset condition
\(\chi(h)^2=\chi(\alpha)^{-1}\).
These degree-one square-map packets have no hidden high-overlap energy.  If
\(M_1,M_2\) are nonconstant degree-one rational parameters and
\(\chi^a(M_1)\chi^b(M_2)\) is geometrically trivial for \(0<a,b<e\), then
either
\[
  M_2=\lambda M_1,\qquad a+b\equiv0\pmod e,
\]
or
\[
  M_2=\lambda/M_1,\qquad a-b\equiv0\pmod e.
\]
Indeed, otherwise one zero or pole appears in only one of the two degree-one
divisors and has nonzero valuation modulo \(e\).  Hence two square-map packets
intersect with the expected product density \(4e^{-2}p+O_e(\sqrt p)\) unless
their parameters are quotient-parallel or quotient-inverse; in the exceptional
cases the intersection is again a single explicit quotient-class packet and is
bounded by \(2|D|+O_e(1)\).
Equivalently, high-overlap square-map packets are grouped by their unordered
zero-pole support \(\Pi=\{\mathrm{zero}(M),\mathrm{pole}(M)\}\subset\mathbb P^1\).
For fixed \(\Pi=\{\alpha,\beta\}\), choose \(M_\Pi\) with divisor
\(\alpha-\beta\).  Every degree-one parameter with this support is
\(\lambda M_\Pi\) or \(\lambda/M_\Pi\), and all its two-coset packets are among
\[
  P_\Pi(\theta)=\{z:\chi(M_\Pi(z))^2=\theta\},
  \qquad \theta\in(\mathbb Z/e\mathbb Z)^2 .
\]
There are exactly \(e/2\) such packets; they are disjoint on
\(\mathbb P^1\setminus\Pi\) and partition that open line.  Thus a
quotient-parallel/inverse high-overlap component contributes only an
\(O_e(1)\) finite palette of distinct slope sets per zero-pole support.
In the square-norm branch the support itself is determined by the norm
boundary.  If \(N_R=B_R/Q_R=\gamma M^2\) with \(M\) nonconstant of degree one,
then
\[
  \operatorname{div}(N_R)=2\operatorname{div}(M).
\]
After reducing \(B_R/Q_R\), the numerator and denominator are scalar multiples
of squares of linear forms on \(\mathbb P^1\), and their roots are exactly the
zero and pole endpoints of \(\Pi(M)\).  Finite zero endpoints are roots of
\(B_R\), finite pole endpoints are roots of \(Q_R\), and infinity accounts for
degree imbalance.  Thus the square-norm palette is indexed by the bounded
norm-boundary divisor of the root-core recurrence chart; changing the square
root only rescales or inverts \(M\), so it stays in the same finite palette.
Since \(B_R\) and \(Q_R\) have degree at most two, this square-norm branch is
also a repeated-endpoint gate.  After cancelling the gcd of \(B_R\) and
\(Q_R\), every finite zero or pole of the reduced norm has valuation
\(\pm2\), and the point at infinity has valuation \(0\) or \(\pm2\).
Equivalently, each nonconstant reduced quadratic factor has discriminant zero
and reduced degree-one factors cannot occur.  Thus the square-norm packet
image lies on the repeated-root/discriminant-zero norm-boundary locus, rather
than on the full generic root image of \(B_R\) and \(Q_R\).
Equivalently, every finite endpoint has a derivative certificate: with
\(\bar B_R,\bar Q_R\) the reduced numerator and denominator, a zero endpoint
satisfies \(\bar B_R(z_0)=\bar B_R'(z_0)=0\) and \(\bar Q_R(z_0)\ne0\), while
a pole endpoint satisfies \(\bar Q_R(z_0)=\bar Q_R'(z_0)=0\) and
\(\bar B_R(z_0)\ne0\).  Hence the exceptional support image is a
double-root/discriminant-zero projection as the root core varies, not a
generic moving-quadratic root image.
For finite square-norm endpoints the bars can be removed: a common finite
factor of \(B_R\) and \(Q_R\) would leave valuation of absolute value at most
one after cancellation, impossible for a square-norm endpoint.  Thus finite
zero endpoints satisfy \(B_R=B_R'=0\) and \(Q_R\ne0\), while finite pole
endpoints satisfy \(Q_R=Q_R'=0\) and \(B_R\ne0\).
Writing \(B_R=b_0+b_1z+b_2z^2\) and \(Q_R=q_0+q_1z+q_2z^2\), this means a
live finite endpoint forces \(b_1^2-4b_0b_2=0\) or \(q_1^2-4q_0q_2=0\); when
the quadratic coefficient is nonzero the endpoint slope is
\(-b_1/(2b_2)\) or \(-q_1/(2q_2)\).  Hence moving square-norm supports are
controlled by raw quadratic discriminants.
In the recurrence rows \(c_i(z)=a_i+zb_i\), these quadratics are the adjacent
Hankel minors
\[
  H_i(z)=c_i(z)c_{i+2}(z)-c_{i+1}(z)^2,
\]
with coefficients
\[
  h_{i,0}=a_i a_{i+2}-a_{i+1}^2,\quad
  h_{i,1}=a_i b_{i+2}+a_{i+2}b_i-2a_{i+1}b_{i+1},\quad
  h_{i,2}=b_i b_{i+2}-b_{i+1}^2.
\]
Thus pole endpoints force the discriminant of \(H_0=Q_R\) to vanish, and zero
endpoints force the discriminant of \(H_1=B_R\) to vanish.
Equivalently, if
\[
 P_i=a_i b_{i+1}-a_{i+1}b_i,\quad
 R_i=a_{i+1}b_{i+2}-a_{i+2}b_{i+1},\quad
 S_i=a_i b_{i+2}-a_{i+2}b_i,
\]
then \(\operatorname{disc}(H_i)=S_i^2-4P_iR_i\).  Endpoint production
therefore forces the adjacent row-pair Plucker minors onto the conic
\(S_i^2=4P_iR_i\), with the proportional-row case as the rank-one
degeneracy.
On the chart \(P_i\ne0\), this conic is parametrized by
\(\lambda_i=S_i/(2P_i)\), so \(S_i=2P_i\lambda_i\) and
\(R_i=P_i\lambda_i^2\).  On the complementary chart \(P_i=0\), one has
\(S_i=0\); if \(R_i\ne0\), then the first row-pair is zero, and symmetrically
if \(R_i=0\) and \(P_i\ne0\), then the third row-pair is zero.  Thus the
endpoint-support ledger splits into zero-row/proportional degeneracies and
the residual moving Plucker chart.
In row form, if \(r_j=(a_j,b_j)\) and \(P_i\ne0\), then
\(r_i,r_{i+1}\) are a basis and the chart is equivalent to
\[
        r_{i+2}=2\lambda_i r_{i+1}-\lambda_i^2 r_i,
        \qquad \lambda_i={S_i\over 2P_i}.
\]
The reverse chart \(R_i\ne0\) gives the symmetric formula
\(r_i=2\mu_i r_{i+1}-\mu_i^2r_{i+2}\), \(\mu_i=S_i/(2R_i)\).  Thus, after the
degenerate charges are removed, the residual endpoint-support row triple is a
one-parameter recurrence rather than a generic conic in row space.
Equivalently, on the \(P_i\ne0\) chart the adjacent Hankel minor itself
factors as
\[
        H_i=c_ic_{i+2}-c_{i+1}^2
          =-(c_{i+1}-\lambda_i c_i)^2 .
\]
On the reverse chart it is
\[
        H_i=-(c_{i+1}-\mu_i c_{i+2})^2 .
\]
Thus the double endpoint is the zero of an explicit row-linear form, with the
constant-linear case accounting for the projective endpoint at infinity.
In particular, when \(b_{i+1}-\lambda_i b_i\ne0\), the finite endpoint is
\[
        z_i={\lambda_i a_i-a_{i+1}\over b_{i+1}-\lambda_i b_i},
\]
and when this denominator vanishes the numerator is nonzero because
\(P_i\ne0\), so the double endpoint is at infinity.  The reverse chart gives
\[
        z_i={\mu_i a_{i+2}-a_{i+1}\over b_{i+1}-\mu_i b_{i+2}} .
\]
Thus a fixed adjacent row basis sends the Plucker chart parameter to endpoint
slope by a projective line map, not by an arbitrary quadratic root map.
The adjacent zero and pole endpoint conditions also overlap rigidly.  If the
two Plucker conics for \(i=0,1\) both hold and \(P_0\ne0\), then with
\(\lambda_0=S_0/(2P_0)\),
\[
        r_2=2\lambda_0r_1-\lambda_0^2r_0,\qquad
        P_1=R_0=P_0\lambda_0^2.
\]
If \(\lambda_0=0\), this is the proportional-row degeneracy for the second
triple.  If \(\lambda_0\ne0\), then \(P_1\ne0\) and the second chart gives
\[
        r_3=2\lambda_1r_2-\lambda_1^2r_1
           =-2\lambda_1\lambda_0^2r_0
            +(4\lambda_0\lambda_1-\lambda_1^2)r_1 .
\]
Thus simultaneous finite endpoint production is either charged already or is
contained in an explicit two-step row-recurrence packet.
In the finite-finite nondegenerate case this packet is recovered from its
ordered endpoint pair.  Since \(P_0\ne0\), \(c_0\) and \(c_1\) have no common
finite zero, so an \(H_0\) endpoint \(z_0\) gives
\[
        \lambda_0={c_1(z_0)\over c_0(z_0)} .
\]
Then an \(H_1\) endpoint \(z_1\) has \(c_1(z_1)\ne0\) and gives
\[
        \lambda_1=2\lambda_0-\lambda_0^2{c_0(z_1)\over c_1(z_1)} .
\]
Thus, for fixed initial row basis \(r_0,r_1\), an ordered finite zero/pole
endpoint pair supports at most one nondegenerate overlapping Plucker packet.
Equivalently, this inversion is projective.  For \(E=[Z:W]\), write
\(c_j(E)=a_jW+b_jZ\).  If \(E_0,E_1\in\mathbb P^1\) are the ordered double
endpoints of \(H_0,H_1\), then \(P_0\ne0\) gives
\[
        \lambda_0={c_1(E_0)\over c_0(E_0)},\qquad
        \lambda_1=2\lambda_0-\lambda_0^2{c_0(E_1)\over c_1(E_1)} .
\]
Thus the projective endpoint pair, including infinity cases, determines the
nondegenerate overlapping Plucker packet over a fixed initial row basis.
If the two projective endpoints coincide, then the inversion formulas give
\(\lambda_1=\lambda_0\).  Consequently
\[
        c_2-\lambda_1c_1
        =\lambda_0(c_1-\lambda_0c_0),
        \qquad H_1=\lambda_0^2H_0 .
\]
The corresponding norm \(H_1/H_0\) is constant, so the diagonal projective
endpoint case is charged to the constant-norm packet ledger rather than to
the residual nonconstant square-norm branch.
After the zero-row/proportional and constant-norm ledgers are removed, the
nondegenerate overlapping chart therefore injects into ordered off-diagonal
projective endpoint pairs.  For an endpoint palette \(\Omega\), a fixed
initial row basis supports at most \(|\Omega|(|\Omega|-1)\) nonconstant
overlapping square-norm packets.
In fact, if \(A_i=\{E:c_i(E)=0\}\), then \(A_0\ne A_1\), the first endpoint
must avoid \(A_0\cup A_1\), and the second must avoid \(A_1\) and the first
endpoint.  Conversely every such ordered pair recovers a unique
nondegenerate packet.  Hence the exact fixed-basis count is
\[
        |\Omega\setminus(A_0\cup A_1)|\,
        (|\Omega\setminus A_1|-1),
\]
which is \((q-1)^2\) for \(\Omega=\mathbb P^1(F_q)\).
In the fixed-basis coordinate \(x=c_1/c_0\), this is the elementary normal
form \(x_0=\lambda_0\in F^\times\),
\[
        x_1={c_1(E_1)\over c_0(E_1)},\qquad
        \lambda_1=2x_0-{x_0^2\over x_1},
\]
with \(x_1\in F^\times\), \(x_1\ne x_0\), and
\[
        H_1=-{x_0^4\over x_1^2}(c_1-x_1c_0)^2 .
\]
The endpoint \(x_1=\infty\) gives \(\lambda_1=2x_0\) and
\(H_1=-x_0^4c_0^2\).
Equivalently the fixed-basis chart is exactly the slope-pair palette
\[
        \lambda_0\in F^\times,\qquad
        \lambda_1\in F\setminus\{\lambda_0\}.
\]
If \(\lambda_1\ne2\lambda_0\), then
\[
        x_1={\lambda_0^2\over 2\lambda_0-\lambda_1},
\]
while \(\lambda_1=2\lambda_0\) is the endpoint \(x_1=\infty\).  Thus
the local nondegenerate branch has no hidden endpoint constraint after the
initial row basis is fixed; the remaining M1 issue is global row-basis
variation.
In these slope-pair coordinates the square map itself is
\[
  m_{\lambda_0,\lambda_1}
  ={(2\lambda_0-\lambda_1)c_1-\lambda_0^2c_0
    \over c_1-\lambda_0 c_0},
  \qquad {H_1\over H_0}=m_{\lambda_0,\lambda_1}^2 .
\]
It is a degree-one bijection from the open projective slope line to
\(F_q^\times\).  Hence the full projective packet has exact multiplicative
character balance, and passing to a finite-slope chart changes a nonprincipal
character sum by at most the one omitted point at infinity.  This fixed-basis
branch has no local \(\sqrt q\) Kummer loss.
The same support view removes orientation duplicates.  In the fixed-basis
coordinate the square-packet palette depends only on the unordered support
\(\{x_0,x_1\}\subset \mathbb P^1(F_q)\setminus\{0\}\), since reversing the
orientation replaces the degree-one parameter by a constant multiple of its
inverse.  Thus the distinct fixed-basis square-packet slope sets are bounded by
\[
        {e\over2}\binom{q}{2}={e\over4}q(q-1),
\]
for even quotient order \(e\).  The ordered endpoint-pair count remains a
packet-incidence ledger, but the support-palette count is the sharper
slope-set ledger.
The slope-pair map to supports is exact: the support
\(\{\lambda_0,x_1\}\) has two preimages when both endpoints are finite
nonzero, corresponding to the two orientations, and one preimage when the
support contains \(\infty\).  Hence
\[
        2\binom{q-1}{2}+(q-1)=(q-1)^2
\]
is precisely the ordered incidence count, while \(\binom q2\) is the support
count.  No further local multiplicity is hidden in the slope-pair chart.
The corresponding support-packet family is an exact incidence design on
\(\mathbb P^1(F_q)\): a point \(x\) lies in exactly \(\binom q2\) packets if
\(x=0\), and exactly \(\binom{q-1}{2}\) packets if
\(x\in\mathbb P^1(F_q)\setminus\{0\}\).  Thus the full local packet family has
no hidden point concentration; nonuniformity must come from the global
row-basis/core-image selection of supports.
More generally, for any selected support family \(\Sigma\) and any selected
packet classes \(\Theta_S\), the packet incidence satisfies
\[
        I(x)\le \#\{S\in\Sigma:x\notin S\},
        \qquad
        \sum_x I(x)={2(q-1)\over e}\sum_{S\in\Sigma}|\Theta_S|.
\]
Equality in the first bound holds for the full class palette.  Thus the
global fixed-basis square-map problem reduces to the support-avoidance profile
of the selected row-basis/core image.
For full class palettes, if \(m=|\Sigma|\) and \(d_\Sigma(x)\) is the degree of
\(x\) in the selected support graph on \(U=\mathbb P^1(F_q)\setminus\{0\}\)
with \(d_\Sigma(0)=0\), then
\[
        I_{\rm full}(x)=m-d_\Sigma(x),\qquad
        \sum_x I_{\rm full}(x)^2=(q-3)m^2+\sum_{x\in U}d_\Sigma(x)^2 .
\]
Thus full-palette concentration is exactly endpoint-degree energy of the
selected support graph.
Equivalently,
\[
        \sum_{x\in U}d_\Sigma(x)^2=2m+2C_1(\Sigma),
\]
where \(C_1(\Sigma)\) is the number of unordered pairs of selected supports
sharing exactly one endpoint.  Thus the excess energy is precisely star-pair
collision energy in the selected support image.
If \(\Delta=\max_x d_\Sigma(x)\), then
\[
        C_1(\Sigma)\le m(\Delta-1),\qquad
        \sum_x I_{\rm full}(x)^2\le(q-3)m^2+2m\Delta .
\]
Consequently a full-palette second moment above
\((q-3)m^2+2mD\) forces an endpoint star of degree \(>D\).
The same bound applies to arbitrary partial class selections, since their
incidence function is pointwise bounded by the full-palette incidence:
\[
        \sum_x I(x)^2\le(q-3)m^2+2m\Delta .
\]
Thus partial square-coset choices do not create a new local source of
second-moment mass.
Writing \(K=\sum_{S\in\Sigma}|\Theta_S|\),
\(M=2(q-1)K/e\), and \(B=|\{x:I(x)>0\}|\), Cauchy's inequality gives
\[
        B\ge {M^2\over (q-3)m^2+2m\Delta}.
\]
More generally, if one tests an endpoint-degree cap \(D\), then packing the
same selected square-map packets into fewer than
\[
        {M^2\over (q-3)m^2+2mD}
\]
slope points forces some endpoint star of degree \(>D\).  Thus unexpectedly
tight packing of many selected square-map packets forces the same
endpoint-star alternative.
This can be turned into a pruning reduction.  Repeatedly remove all supports
incident to an endpoint whose current degree is \(>D\).  The deleted supports
are partitioned among endpoint-star templates, while the residual support
family \(\Sigma_{\rm ap}\) has maximum degree at most \(D\).  For the residual
classes, with
\[
        K_{\rm ap}=\sum_{S\in\Sigma_{\rm ap}}|\Theta_S|,\quad
        M_{\rm ap}=2(q-1)K_{\rm ap}/e,\quad
        B_{\rm ap}=|\{x:I_{\rm ap}(x)>0\}|,
\]
one has
\[
        B_{\rm ap}\ge
        {M_{\rm ap}^2\over (q-3)m_{\rm ap}^2+2m_{\rm ap}D}
        \qquad (m_{\rm ap}=|\Sigma_{\rm ap}|>0).
\]
Thus after endpoint stars are charged to quotient/low-template ledgers, the
aperiodic residual cannot be overpacked by local square-map geometry.
Equivalently, for any residual slope budget \(R\), either \(B_{\rm ap}>R\) or
\[
        \left({2(q-1)\over e}\right)^2 K_{\rm ap}^2
        \le R\bigl((q-3)m_{\rm ap}^2+2m_{\rm ap}D\bigr),
\]
or, in average-density form,
\[
        (K_{\rm ap}/m_{\rm ap})^2
        \le {R((q-3)+2D/m_{\rm ap})\over (2(q-1)/e)^2}.
\]
Thus a dense bounded-degree residual must occupy many slopes; if it does not,
the residual square-coset class selection is sparse rather than geometrically
overpacked.
In the full-palette case there is an exact sharper support formula:
\[
        \{x:I_{\rm full}(x)>0\}
        =\mathbb P^1(F_q)\setminus\bigcap_{S\in\Sigma}S,\qquad
        B_{\rm full}=q+1-\left|\bigcap_{S\in\Sigma}S\right|.
\]
Hence, after support-star pruning, if \(m_{\rm ap}>D\), the residual common
intersection is empty and the full-palette residual covers all \(q+1\)
projective slopes.  Any small pruned residual support therefore comes from at
most \(D\) remaining supports or from genuinely partial square-coset palettes.
The partial-palette leftover has an exact defect ledger.  If
\(\Theta_S^c\) denotes the missing square-coset classes and
\[
        J(x)=\#\{(S,\theta):\theta\in\Theta_S^c,\ x\in P_S(\theta)\},
\]
then
\[
        I(x)+J(x)=m-d_\Sigma(x),\qquad
        \sum_x J(x)={2(q-1)\over e}\sum_{S\in\Sigma}|\Theta_S^c|.
\]
Thus, after pruning with maximum residual degree \(D\) and \(m_{\rm ap}>D\),
the number of uncovered residual slopes satisfies
\[
        |\{x:I_{\rm ap}(x)=0\}|
        \le {2(q-1)\over e}\,
        {\sum_{S\in\Sigma_{\rm ap}}|\Theta_S^c|\over m_{\rm ap}-D}.
\]
So a large bounded-degree residual can leave many slopes uncovered only by
spending many missing local square-coset classes.
Combining the two residual ledgers gives the local support-budget alternative.
For any \(0\le R\le q+1\), after pruning endpoint stars with degree cap \(D\),
either \(m_{\rm ap}\le D\), or \(B_{\rm ap}>R\), or both
\[
        \left({2(q-1)\over e}\right)^2 K_{\rm ap}^2
        \le R\bigl((q-3)m_{\rm ap}^2+2m_{\rm ap}D\bigr)
\]
and
\[
        (m_{\rm ap}-D)(q+1-R)
        \le {2(q-1)\over e}\, C_{\rm ap}
\]
hold.  Thus small residual slope support must be paid for by both selected
class sparsity and missing-class mass.
Setting \(R=B_{\rm ap}\) gives a direct palette-density lower bound.  With
\(h=e/2\),
\[
        \alpha_{\rm ap}=K_{\rm ap}/(hm_{\rm ap}),\qquad
        \gamma_{\rm ap}=C_{\rm ap}/(hm_{\rm ap})=1-\alpha_{\rm ap},
\]
and \(m_{\rm ap}>D\),
\[
        B_{\rm ap}\ge
        {(q-1)^2\alpha_{\rm ap}^2m_{\rm ap}^2
         \over (q-3)m_{\rm ap}^2+2m_{\rm ap}D},
        \qquad
        B_{\rm ap}\ge
        q+1-{(q-1)\gamma_{\rm ap}m_{\rm ap}\over m_{\rm ap}-D}.
\]
In particular, if \(m_{\rm ap}\ge2D\) and
\(\gamma_{\rm ap}\le\varepsilon\), then
\(B_{\rm ap}\ge q+1-2\varepsilon(q-1)\).
Equivalently, for a target support budget \(R\), if \(m_{\rm ap}>D\) and
either
\[
        \left({2(q-1)\over e}\right)^2K_{\rm ap}^2
        > R\bigl((q-3)m_{\rm ap}^2+2m_{\rm ap}D\bigr)
\]
or
\[
        {2(q-1)\over e}C_{\rm ap}
        <(m_{\rm ap}-D)(q+1-R),
\]
then \(B_{\rm ap}>R\).  Thus a global proof can close this local branch by
showing enough selected palette density or too few missing palette classes.
The local fixed-basis square-map output is therefore a four-way theorem: after
greedy endpoint-star pruning at cap \(D\), the deleted part is charged to
endpoint-star templates; or \(m_{\rm ap}\le D\); or \(B_{\rm ap}>R\); or the
small-support residual satisfies both support-budget inequalities above.  The
last case is excluded by either displayed strict density trigger, so the only
remaining small-support branch is a genuine global palette-sparsity
certificate.
Equivalently, the optional certificate is the exact finite object
\[
        m_{\rm ap}>D,\quad B_{\rm ap}\le R,\quad
        \max_x d_{\Sigma_{\rm ap}}(x)\le D,
\]
together with the two support-budget inequalities.  Thus the next global M1
task is to rule out this sparse certificate from the row-basis/core image, or
to identify it as quotient/template structure.
With \(h=e/2\), \(\alpha=K/(hm)\), and
\(\gamma=C/(hm)=1-\alpha\), any such certificate lies in the feasibility
window
\[
        \alpha^2\le
        {R((q-3)m^2+2mD)\over (q-1)^2m^2},
        \qquad
        \gamma\ge{(m-D)(q+1-R)\over(q-1)m}.
\]
So a global lower bound on selected palette density above this window rules
out the sparse certificate in the corresponding parameter range.
If additionally \(m\ge LD\) for an integer \(L\ge2\), this becomes the uniform
far-from-star window
\[
        \alpha^2\le {R(q-3+2/L)\over(q-1)^2},
        \qquad
        \gamma\ge {L-1\over L}\,{q+1-R\over q-1}.
\]
Thus certificates whose residual size is a fixed factor beyond the endpoint
degree cap can be excluded by a single \(L\)-dependent density threshold.
Contrapositively, if either strict density gap violates this window, then
\[
        m<LD.
\]
The residual then has fewer than \(LD\) two-point supports and endpoint
footprint \(<2LD\).  Hence a global row-basis/core-image density theorem
above the displayed far-from-star window reduces every remaining
small-support residual to a bounded near-star/template ledger.
Writing \(h=e/2\) and
\[
        s_{LD}=\min(q+1,\max(0,2LD-1)),
\]
the number of such residual palette templates is bounded by
\[
        \sum_{s=0}^{s_{LD}} {q+1\choose s}\,2^{h{s\choose 2}}
        \le (s_{LD}+1)(q+1)^{s_{LD}}2^{h{s_{LD}\choose 2}}.
\]
For fixed \(D,L,e\), the near-star branch is therefore a polynomial-size
endpoint-template ledger.
Equivalently, define
\[
        \theta_L(q,R)
        =\min\left(
          {\sqrt{R(q-3+2/L)}\over q-1},
          1-{L-1\over L}{q+1-R\over q-1}
        \right).
\]
If a global row-basis/core-image argument proves
\(\alpha_{\rm ap}>\theta_L(q,R)\) for every sparse residual output, then the
far-from-star sparse branch is empty.  The local output reduces to endpoint
stars, \(D\)-small residuals, large support \(B_{\rm ap}>R\), and the
polynomial-size near-star template ledger above.
Equivalently, a density lower bound \(\alpha_{\rm ap}\ge a_0\) gives the
explicit support budget
\[
        R_{\rm dens}(a_0,L)=
        \max\left\{
          {a_0^2(q-1)^2\over q-3+2/L},
          q+1-{L\over L-1}(1-a_0)(q-1)
        \right\}.
\]
For every \(R<R_{\rm dens}(a_0,L)\), a sparse certificate with
\(B_{\rm ap}\le R\) is either \(L\)-near-star or impossible.  Thus a global
selected-density theorem immediately translates into an explicit
all-line support budget plus the near-star template ledger.
Equivalently, if \(D\ge1\) and an endpoint-footprint cap \(S\ge4D-1\) is
chosen, set
\[
        L_S=\left\lfloor {S+1\over 2D}\right\rfloor .
\]
The budget \(R_{\rm dens}(a_0,L)\) is nondecreasing in \(L\), so \(L_S\) is
the largest far-factor compatible with footprint at most \(S\).  Under
\(\alpha_{\rm ap}\ge a_0\), every sparse certificate with
\(B_{\rm ap}<R_{\rm dens}(a_0,L_S)\) is impossible or has endpoint footprint
at most \(S\), leaving at most
\[
        \sum_{s=0}^{\min(q+1,S)} {q+1\choose s}\,2^{h{s\choose2}}
\]
palette templates.
For finite rows one can replace the real comparison by the integer ceiling
\[
        R_{\mathbb Z}(a_0,L)=\lceil R_{\rm dens}(a_0,L)\rceil-1.
\]
Then an integer support budget satisfies
\(R<R_{\rm dens}(a_0,L)\) exactly when \(R\le R_{\mathbb Z}(a_0,L)\).
Thus, with \(L=L_S\), every sparse certificate with
\(B_{\rm ap}\le R_{\mathbb Z}(a_0,L_S)\) is impossible or belongs to the
endpoint-footprint-\(\le S\) template ledger.
There is a weak unconditional floor in the reduced support convention:
each residual support carries at least one selected square-coset class.  Hence
\[
        \alpha_{\rm ap}\ge {1\over h}={2\over e}.
\]
Substituting \(a_0=1/h\) in \(R_{\rm dens}\) gives a theorem-backed baseline
budget \(R_{\rm base}(L)\); every sparse certificate below this budget is
already impossible or \(L\)-near-star before any stronger global density
input is used.
Moreover, if \(L_E(z)=Wz-Z\) for \(E=[Z:W]\), then in the off-diagonal case
\[
        {H_1\over H_0}
        =\gamma^2\left({L_{E_1}\over L_{E_0}}\right)^2
\]
for some \(\gamma\in F^\times\).  Thus this residual chart is exactly a
degree-one square-map packet with divisor \(2E_1-2E_0\), up to an overall
square scalar.
For even character order \(e\), fixing the ordered pair \(E_0,E_1\) leaves
only the \(e/2\)-packet square-coset palette
\[
        P_{E_0,E_1}(\theta)
        =\{z:\chi(L_{E_1}(z)/L_{E_0}(z))^2=\theta\},
        \qquad \theta\in 2\mathbb Z/e\mathbb Z,
\]
which partitions the open slope line away from \(E_0,E_1\).
Combining the endpoint-pair injection and this palette, a fixed initial row
basis and endpoint palette \(\Omega\) contribute at most
\((e/2)|\Omega|(|\Omega|-1)\) distinct nonconstant slope sets; for the full
projective line over \(F_q\), this is \((e/2)q(q+1)\).
Each such packet has at most \(2(q-1)/e\) finite slopes.  If both endpoints
are finite, exactly one packet loses the point at infinity and has size
\(2(q-1)/e-1\); otherwise all packets have the full \(2(q-1)/e\) finite
slopes.
The finite endpoints are also ledger charges rather than residual live slopes.
If \(B_R(z_0)=0\) and \(Q_R(z_0)\ne0\), then
\(p_R(z_0)=B_R(z_0)/Q_R(z_0)=0\), so the recovered locator lies on the
fixed-root line \(p=0\).  If \(Q_R(z_0)=0\), then \(z_0\) is outside the
recurrence chart; whenever the recurrence is solvable, the denominator-zero
classification gives the existing full-plane or fixed-root charge, and if it
is not solvable there is no formal two-root extension.  The endpoint at
infinity is the same statement in the reciprocal slope chart.
Consequently, for any finite family \(\mathcal S\) of degree-one square-map
packets, the number of distinct slope sets in \(\mathcal S\) is at most
\[
        \frac e2\,
        \#\{\Pi(M):P(M,\kappa)\in\mathcal S\}.
\]
In particular a single nonconstant square-norm branch contributes at most
\(e/2\) distinct square-map slope sets; repeated square roots and
parallel/inverse parameters are multiplicity inside the same endpoint palette,
not a new slope-growth source.
This is cruder than the cover bound below, but it is genus-free; the cover
sums are only needed for the extra factor-\(e\) saving.  Now define
\[
  S_a^+=\sum_{C_R^\times}\chi^a(r_+),\qquad
  S_{a,b}=\sum_{C_R^\times}\chi^a(r_+)\chi^b(r_-).
\]
The ordered mixed-domain count has the exact expansion
\[
  N_R^{+,-}(D)=
  \frac1e\sum_a S_a^+ - \frac1{e^2}\sum_{a,b}S_{a,b}+O(1).
\]
Its principal contribution is \(((e-1)/e^2)|C_R^\times|+O(1)\).  If no
nonprincipal term is geometrically trivial, the standard curve Kummer-Weil
input gives
\[
  N_R^{+,-}(D)
  \le \frac{e-1}{e^2}|C_R^\times|+C_e\sqrt p+O_e(1).
\]
The norm map gives a necessary condition for any cover-level power branch:
if \(F_{a,b}=r_+^a r_-^b\) has divisor divisible by \(e\) on the cover, then
pushing forward by \(C_R\to\mathbb P^1_z\) gives
\[
  \mathrm{div}\bigl((B_R/Q_R)^{a+b}\bigr)\equiv0\pmod e.
\]
Hence, after the constant-norm and square-norm branches above are charged or
isolated, all one-root terms and all pair terms with \(a+b\not\equiv0\pmod e\)
are nontrivial bounded-conductor cover sums.  The only remaining cover-level
power target is the anti-diagonal ratio term \(\chi^a(r_+/r_-)\).
This target is almost closed as well.  Let \(\Phi_R=r_+/r_-\).  Since
\((A_R+Y)(A_R-Y)=4Q_RB_R\), the map \(\Phi_R:C_R\to\mathbb P^1\) has degree
at most four, and
\[
  \Phi_R/N_R=r_-^{-2},\qquad N_R=B_R/Q_R.
\]
Thus even-order anti-ratio power branches are already square-norm branches,
while odd orders \(m\ge5\) are impossible for nonconstant \(\Phi_R\).  After
the square-norm branch is isolated, the only possible nonconstant anti-ratio
power branch is the cubic one \(\mathrm{div}_C(r_+/r_-)\equiv0\pmod3\).
In particular, if \(3\nmid e\), no cover-level power branch remains after the
norm exceptions.
If \(3\mid e\), geometric triviality of this cubic term would require
\(\Phi_R=cG^3\) in \(k(C_R)^*\).  Since \(\deg\Phi_R\le4\), nonconstant
\(\Phi_R\) then has \(\deg\Phi_R=3\) and \(\deg G=1\).  Thus a genuine cubic
no-cancellation term forces the normalization of \(C_R\) to be rational; on a
geometrically integral genus-one cover the cubic anti-ratio term is still
geometrically nontrivial and contributes only a bounded-conductor Weil term.
On the rational cubic branch, the only possible large anti-diagonal
coefficients are \(a=e/3,2e/3\), when \(3\mid e\).  Hence, after the
constant-norm and square-norm exceptions are removed, their conjugate cubic
constants contribute at worst \(+|C_R^\times|/e^2\) in the mixed-domain
expansion, and
\[
  N_R^{+,-}(D)
  \le {1\over e}|C_R^\times|+C_e\sqrt p+O_e(1)
  \le |D|+C_e\sqrt p+O_e(1)
\]
for geometrically integral rational normalization.
The diagonal pair terms in this expansion descend to the slope line:
\[
  S_{a,a}
  =\sum_z (1+\chi_2(\Theta_R(z)))\chi^a(B_R(z)/Q_R(z)),
\]
on the open locus \(Q_RB_R\ne0\).  Hence each diagonal term is a sum of two
bounded-support \(\mathbb P^1\) Kummer traces, with support contained in the
roots of \(B_R,Q_R,\Theta_R\) and infinity.  Its only no-cancellation branches
are the explicit power-divisor congruences for \((B_R/Q_R)^a\) and
\((B_R/Q_R)^{a\ell/e}\Theta_R^{\ell/2}\), \(\ell=\mathrm{lcm}(e,2)\).
Only the off-diagonal \(S_{a,b}\), \(a\ne b\), and the one-root sums
\(S_a^+\) remain genuinely cover-level terms.
For index two, these cover-level terms cancel.  Writing \(\rho=\chi\), sheet
symmetry gives \(S_1^+=S_1^-\), hence
\[
  N_R^{+,-}(D)
  =\frac14\left(|C_R^\times|-S_{1,1}\right)+O(1)
  =\frac14|C_R^\times|
   -\frac14\sum_z(1+\rho(\Theta_R(z)))\rho(B_R(z)/Q_R(z))+O(1).
\]
Thus quadratic-residue domains reduce the mixed root-core count entirely to
bounded-support \(\mathbb P^1\) Kummer sums after the fixed-zero-root branch is
charged.  If neither \(B_R/Q_R\) nor \(\Theta_RB_R/Q_R\) is a square divisor,
the nonsplit branch is at most \(|D|/2+C\sqrt p+O(1)\), and the split-square
branch is at most \(|D|+C\sqrt p+O(1)\).
Thus a geometrically integral nonsplit quartic-cover branch has at most
\((1-1/e)|D|+C_e\sqrt p+O_e(1)\) points, with conjugate or no-point
degeneracies only smaller, while the split-square rational branch has at most
\(2(1-1/e)|D|+C_e\sqrt p+O_e(1)\).  Each unordered mixed split pair
contributes to exactly one ordered sheet, and repeated roots contribute
nothing to the mixed-domain condition.
Combining the cover-power reductions gives the classified per-core form:
after fixed-zero-root and constant-norm branches are charged and the
negative-parity rational one-root-square subbranch is isolated,
positive-parity square norm satisfies \((e-2)|D|+C_e\sqrt p+O_e(1)\),
geometrically integral genus-one
branches satisfy \((1-1/e)|D|+C_e\sqrt p+O_e(1)\), negative-parity
square-norm genus-one branches satisfy \(|D|+C_e\sqrt p+O_e(1)\),
geometrically integral rational cubic branches satisfy
\(|D|+C_e\sqrt p+O_e(1)\), rational one-root-square branches satisfy
\(O_e(\sqrt p)\) for the negative root-square sign and
\(2|D|+C_e\sqrt p+O_e(1)\) for the positive root-square sign, and
split-square rational branches satisfy
\(2(1-1/e)|D|+C_e\sqrt p+O_e(1)\).  The index-two nonsplit case away from the
two square gates has the sharper \(|D|/2+C\sqrt p+O(1)\) bound above.

Consequently, after fixed-slope root slices are removed, the residual \(t=2\)
one-exchange graph has maximum degree at most \(j\): each locator has only
\(j\) possible \((j-1)\)-cores, and each core supplies at most one residual
neighbor.  In particular the unordered residual edge count is at most
\(j|A|/2\), and also at most \(\binom{|D|}{j-1}\) by counting cores.
Plugging this into the average support-collinearity ledger
\path{experimental/notes/m1/m1_average_support_collinearity.md}, and writing
the line-field size as \(Q\), gives for the residual family \(A_{\rm res}\),
\(M=|A_{\rm res}|\),
\[
  B_2^{\max}(A_{\rm res})
  \le
  \frac{1-p_z}{M p_z}+\frac{4jQ}{M},
  \qquad
  p_z=Q^{-2}(1-Q^{-2}).
\]
Thus this branch contributes \(o(1)\) to the average missing-slope density
whenever \(M p_z\to\infty\) and \(M/(jQ)\to\infty\).  The remaining M1
obstruction must come from small residual family size, root-slice charges, or
higher packet/two-exchange structure; this is not yet a worst-case packing
theorem.

The same substitution has a packet-level higher-exchange form.  Inside one
same-slope affine \(h\)-exchange packet, after full moving \(r\)-root fibers
have been charged, the residual codegree bound
\(\Gamma_r\le\binom{h}{r}(Q_F^{r-1}-1)\) gives, for average-ledger support
slack \(\tau\),
\[
  B_\tau^{\max}(A_{\rm res})
  \le
  \frac{1-p_z}{M p_z}
  +\frac{4}{M}\sum_{r=1}^{\min(h,\tau-1)}
    \binom{h}{r}(Q_F^{r-1}-1)Q_F^{\tau-r},
  \qquad
  p_z=Q_F^{-\tau}(1-Q_F^{-\tau}).
\]
This is still a local packet ledger input; a global M1 estimate must also
sum over packet choices and the quotient, tangent/contained, split-root, and
fixed-slope ledgers.

\subsection{M1: variable-line packet reduction}

The follow-up variable-line packet note isolates a local reduction inside the
same Hankel-pencil branch.  Fix a \((j-2)\)-core \(R\) and a non-fixed proper
one-dimensional determinant component \(L\) in the two-root plane
\[
  T=R\cup\{x,y\},\qquad s=x+y,\quad p=xy .
\]
The constant-line collapse above verifies the needed variable-slope hypothesis
for every surviving fixed-sum or nondegenerate product-Mobius packet after
fixed-root and full-plane charges.  Thus the active-new mass \(r_L\) on such
an \(L\) obeys
\[
  r_L
  \le
  {\bf 1}_{d_L=m_L=r_L=1}
  +(d_L-m_L)+2\binom{m_L}{2}.
\]
The unordered packet pairs in the final term inject into the global
different-slope two-exchange codegree ledger, so summing over non-fixed
line packets gives
\[
  \sum_L r_L
  \le S_{\rm dom}+Q_{\rm def}+2E_{\rm pkt}.
\]
The remaining domain-singleton term is reduced to contained/tangent and
one-outside target images:
\[
  S_{\rm dom}
  \le
  Z_0
  +2\binom{j}{2}\binom{n-j}{2}|{\rm Target}_{\rm cb}|
  +(j-1)\binom{n-j+1}{2}|{\rm Target}_{\rm off}|.
\]
Writing \(a=n-j\), the exceptional zero-lower term vanishes whenever
\[
  a>\frac{n+1}{2},
\]
and hence throughout the positive-slack rate-half range \(k\ge n/2\),
\(a=k+t\), \(t\ge1\).  Consequently the non-fixed variable-line branch is
reduced to two explicit residual estimates: the active different-slope
two-exchange codegree and the one-outside boundary shadow image, after quotient
defects, contained/tangent ledgers, and fixed-slope boundary slices have been
paid.  Indeed the shadow-fiber bound gives
\[
  |{\rm Target}_{\rm off}^{\rm res}|
  \le 2|{\rm Shadow}_{\rm off}^{\rm res}|.
\]
After the fixed-anchor boundary-core slices are also charged, the sharper
fiber bound
\[
  (j-1)|{\rm Target}_{\rm off}^{\rm res}|
  \le 2|{\rm Core}_{\rm off}^{\rm res}|
\]
may be substituted directly into the singleton target term.  The boundary
contribution
\[
  (j-1)\binom{n-j+1}{2}|{\rm Target}_{\rm off}^{\rm res}|
\]
is then bounded by
\[
  2\binom{n-j+1}{2}|{\rm Core}_{\rm off}^{\rm res}|,
\]
so a polynomial bound for the boundary-core image closes this branch with only
an \(n^2\) bookkeeping loss, rather than the older \(n^3\) boundary-image loss.
The fixed-core graph form also gives
\[
  |{\rm Core}_{\rm off}^{\rm res}|
  \le 2(n-j+2)|{\rm Root}_{\rm off}^{\rm res}|,
\]
so bounding the active \((j-2)\)-core image closes the same branch with an
\(n^3\) boundary loss after all boundary-core root slices have been charged.

\textbf{Status.}  This is a proved local packet lemma under the Hankel-pencil
normal form after the same-slope, fixed-root, full-plane, and boundary-slice
charges recorded above.  It is not the final all-line M1 theorem and does not
add a leaderboard row.

\subsection{M1: width-one fixed-root closure}

PR \#136, due to AllenGrahamHart, extracts a compact closure criterion for the
width-one critical-tail branch.  At a base-free canonical \(b=2\) node, let
\[
  V=\langle P,Q\rangle\subset F[X]_{<q},\qquad |D^{can}|=q+s.
\]
A root-free width-one certificate is a direction \(A\in V\) whose root shadow
on \(D^{can}\) has maximal size \(q-1\).  Writing \(Z=Z(A)\) and
\(O=D^{can}\setminus Z\), one has
\[
  A=c\ell_Z,\qquad |Z|=q-1,\qquad |O|=s+1.
\]
Equivalently, \(O\) supports a width-one certificate if and only if
\[
  L_O=\ell_{D^{can}\setminus O}\in V,
\]
or equivalently all \(3\times3\) coefficient minors of \(P,Q,L_O\) vanish.
Thus width-one certificates are bounded-complement rank tests rather than an
unstructured projective-direction family.

The large width-one cube descends losslessly under fixed-root absorption.  If
\(P_m\subset Z\), then
\[
  A^{P_m}=A/\ell_{P_m}=c\ell_{Z\setminus P_m},\qquad
  D^{can}\setminus P_m=(Z\setminus P_m)\sqcup O,
\]
so the descended width remains one and the bounded complement \(O\) is
unchanged.  With \(a=\lfloor(q-2)/2\rfloor\), the co-small flag cube has the
first-root partition
\[
  \sum_{e=1}^a\binom{q-1}{e}
  =
  \sum_{x\in Z}\sum_{f=0}^{a-1}\binom{|Z_{>x}|}{f}.
\]
After exposing the first root \(x\), every nonempty width-one flag injects into
a one-root absorbed fixed-divisor/root-slice family.  Quantitatively, for
\[
  M_x=\#\{F\subset Z_{>x}: |F|\le a-1\},
\]
one gets
\[
  \max_{x\in Z}M_x
  \ge {1\over q-1}\sum_{e=1}^a\binom{q-1}{e}
  \ge {2^{q-1}\over 2q^2}\qquad(q\ge4).
\]
Thus a large width-one cube already produces a large one-root slice; it cannot
hide as a many-first-root averaging phenomenon.

Consequently, in fixed surplus \(s_0\le\sigma\), small nodes are polynomial by
the active-tree bound, and on \(q_S>s_S+2\) there is at most one width-one
shadow per node.  The local closure criterion is
\[
  WO_1(A_0)
  \le FixedRootOneRoot_{r1}(A_0)+O_\sigma(Q^{\sigma+1}),
\]
where \(Q\) is the initial quotient width.  This closes the width-one
critical-tail branch only after a polynomial bound for the displayed one-root
fixed-root/root-slice ledger is proved or imported.  It is not the final
all-line M1 theorem and does not add a leaderboard row.

\subsection{Challenge-map pullbacks and public status corrections}

The adjacent-ledger update adds the missing pullback rule from slope ledgers to
protocol challenge maps.  For a challenge-to-slope map \(\phi:K\to Z\), a bad
set \(B\subseteq Z\) pulls back with probability
\[
  \frac{|\phi^{-1}(B)|}{|K|}.
\]
If the nonempty fibers have sizes \(s_1\ge s_2\ge\cdots\), the adversarial
envelope for \(|B|=L\) is \((s_1+\cdots+s_L)/|K|\).  For an \(\F_q\)-linear map
of rank \(r\), the effective denominator is \(q^r\), not the ambient challenge
space size.  This prevents accidental multiplication of denominators in protocol
ledgers.

The public-board status corrections are:
\[
\begin{array}{c|c|c|c}
\text{row} & a & \text{tangent floor} & \text{status}\\
\hline
\text{Cycle116 gate} & 262 & 251 & \text{gate unconditional; exact count conditional}\\
\text{Cycle119 gate} & 263 & 250 & \text{gate unconditional; exact count conditional}\\
\text{reserve272} & 272 & 241 & \text{proved but subsumed}\\
\text{reserve288} & 288 & 225 & \text{proved but subsumed}\\
\text{reserve313} & 313 & 200 & \text{proved but subsumed}
\end{array}
\]
The exact Cycle84 numerator \(52{,}747{,}567{,}092\) remains a
computation-dependent record; only the weaker \(>2^{-128}\) gate follows from
the generic tangent floor without that finite census.

% High-agreement material promoted to clean notes.
\input{notes/high_agreement/tangent_staircase}
\input{notes/high_agreement/line_ca_projective}
\input{notes/high_agreement/curve_mca_ca}
\input{notes/high_agreement/interleaved_uniqueness}
\input{notes/high_agreement/current_row_protocol_ledger}
\input{notes/high_agreement/general_threshold_compiler}
\section{Towards-Prize Finite-Threshold Theorems}
\label{sec:towards-prize-thresholds}

\textbf{Status: \textsc{New-local} finite theorem infrastructure plus
subsumed strict264 proof record.}
The results in this section are deliberately finite and certificate-facing.  They
are meant to move the experimental ledger from lower-bound numerators toward the
integer agreement staircase needed by the Proximity Prize.  They do not prove the
corrected Paper B local-limit conjecture.  The later
high-agreement tangent staircase supersedes the old strict264/265 search target
as a threshold-pinning route for the \(\F_{17^{32}}\), \(n=512\), \(k=256\)
row; the strict264 material below remains useful as certificate language and
as a record of quotient-core mechanisms.

For a received line \(\ell_z=f+zg\), write
\[
  \operatorname{Bad}_{f,g}(C,a)
\]
for the set of finite slopes \(z\in F\) for which there is a support
\(S\subseteq D\), \(|S|\ge a\), and a codeword \(c\in C\) satisfying
\(c|_S=(f+zg)|_S\), but there is no pair of codewords
\(c_f,c_g\in C\) with \(c_f|_S=f|_S\) and \(c_g|_S=g|_S\).  Put
\[
  \operatorname{LD}_{\rm sw}^{f,g}(C,a)
  =|\operatorname{Bad}_{f,g}(C,a)|,
  \qquad
  \operatorname{LD}_{\rm sw}(C,a)
  =\max_{f,g}\operatorname{LD}_{\rm sw}^{f,g}(C,a),
\]
where the maximum is over the line family allowed by the sampler under
consideration.  The field paying the denominator is always \(\qline\), not
silently \(\qgen\) or \(\qchal\).

\begin{proposition}[certificate-to-\(\operatorname{LD}_{\rm sw}\) gate]
\label{prop:certificate-to-ldsw}
Let \(C=\RS[F,D,k]\), let \(a\ge k\), and fix a received line \(f+zg\).  Suppose
there is a finite set \(\mathcal C\) of tuples
\[
  (z_i,S_i,P_i),
  \qquad z_i\in F,
  \quad S_i\subseteq D,
  \quad |S_i|\ge a,
  \quad P_i\in F[X]_{<k},
\]
with the following properties:
\begin{enumerate}
\item the slopes \(z_i\) are pairwise distinct;
\item \(P_i(x)=f(x)+z_i g(x)\) for every \(x\in S_i\);
\item there are no \(P_{i,f},P_{i,g}\in F[X]_{<k}\) satisfying
  \(P_{i,f}|_{S_i}=f|_{S_i}\) and \(P_{i,g}|_{S_i}=g|_{S_i}\).
\end{enumerate}
Then
\[
  \operatorname{LD}_{\rm sw}^{f,g}(C,a)\ge |\mathcal C|,
  \qquad
  \operatorname{LD}_{\rm sw}(C,a)\ge |\mathcal C|.
\]
\end{proposition}

\begin{proof}
Each tuple certifies that the slope \(z_i\) is explained on a support of size at
least \(a\), and the third condition is exactly support-wise noncontainment on
that support.  Pairwise distinctness of the slopes prevents multiple tuples from
paying for the same bad slope.  Thus all \(z_i\)'s lie in
\(\operatorname{Bad}_{f,g}(C,a)\), giving the first inequality.  The second
follows from the definition of \(\operatorname{LD}_{\rm sw}(C,a)\) as a maximum
over lines.
\end{proof}

\begin{corollary}[strict264 certificate target]
\label{cor:strict264-certificate-target}
For the row
\[
  C=\RS[\F_{17^{32}},H,256],
  \qquad |H|=512,
\]
a strict264 certificate consists of seven tuples in
Proposition~\ref{prop:certificate-to-ldsw} with \(a=264\).  Such a certificate
proves
\[
  \operatorname{LD}_{\rm sw}(C,264)\ge 7.
\]
\end{corollary}

\begin{proof}
Apply Proposition~\ref{prop:certificate-to-ldsw} with
\(|\mathcal C|=7\) and \(a=264\).
\end{proof}

\begin{proposition}[fixed-locator unique-slope principle]
\label{prop:fixed-locator-unique-slope}
Use the notation of Theorem~\ref{thm:syndrome-pencil}.  Fix
\(a\ge k\), put \(t=a-k\), \(j=n-a=r-t\), and fix a squarefree
\(D\)-split locator \(L_T\) of degree \(j\).  Let
\[
  A_T=H_{t,j}(u)\boldsymbol\ell_T,
  \qquad
  B_T=H_{t,j}(v)\boldsymbol\ell_T.
\]
If \(B_T\ne0\), then this locator contributes at most one finite
support-wise noncontained slope.  More precisely, a slope exists exactly when
\(A_T\) is a scalar multiple of \(B_T\), in which case the unique slope is
\[
  z_T=-A_{T,m}/B_{T,m}
\]
for any coordinate \(m\) with \(B_{T,m}\ne0\), and the same value must be
obtained from every nonzero coordinate of \(B_T\).
Consequently
\[
  \operatorname{LD}_{\rm sw}^{f,g}(C,a)
  \le
  \#\{T\subset D:
        |T|=n-a,
        L_T\text{ squarefree and }D\text{-split},
        B_T\ne0,
        A_T\in F B_T\}.
\]
\end{proposition}

\begin{proof}
For the fixed locator, Theorem~\ref{thm:syndrome-pencil} gives the incidence
equation
\[
  A_T+zB_T=0
\]
and the noncontainment condition \(B_T\ne0\).  If \(B_T\ne0\), any solution
\(z\) is forced by any coordinate where \(B_T\) is nonzero.  Existence is exactly
the consistency condition that all coordinates give the same scalar, equivalently
\(A_T\in F B_T\).  Counting possible locators gives the final bound; different
locators may still collide to the same slope, so the displayed quantity is an
upper bound on distinct slopes.
\end{proof}

\begin{remark}[why this helps the 265 upper bound]
For the current \(n=512,k=256\) row, this proposition remains useful for
duplicate-aware locator searches.  The old proposed upper bound
\(\operatorname{LD}_{\rm sw}(C,265)\le6\) is now known to be false by
Theorem~\ref{thm:moving-root-tangent-floor}, which already gives
\(\operatorname{LD}_{\rm sw}(C,265)\ge248\).  Future upper-bound searches should
therefore start beyond the exact tangent gate, at agreement \(507\) and higher,
or explicitly isolate non-tangent mechanisms.
\end{remark}

\begin{theorem}[subfield confinement for base-valued lines]
\label{thm:subfield-confinement-base-lines}
Let \(K\subseteq F\) be a subfield, let \(D\subset K\), and let
\(C_F=\RS[F,D,k]\).  Suppose \(f,g:D\to K\), and let
\(S\subseteq D\) have \(|S|\ge k\).  If \(z\in F\setminus K\) and there is a
polynomial \(P\in F[X]_{<k}\) with
\[
  P(x)=f(x)+zg(x)
  \qquad (x\in S),
\]
then there are polynomials \(P_f,P_g\in K[X]_{<k}\) such that
\[
  P_f|_S=f|_S,
  \qquad
  P_g|_S=g|_S.
\]
Consequently, for \(K\)-valued received lines and agreement \(a\ge k\), every
support-wise noncontained finite slope lies in \(K\).
\end{theorem}

\begin{proof}
Choose a subset \(S_0\subseteq S\) of size \(k\).  Interpolate the values of
\(f\) and \(g\) on \(S_0\) to obtain unique polynomials
\(P_f,P_g\in K[X]_{<k}\).  On \(S_0\), the polynomial
\(P_f+zP_g\) agrees with \(P\).  Since both have degree less than \(k\) and agree
on \(k\) points, they are equal as polynomials over \(F\).  For every
\(x\in S\), therefore,
\[
  0=f(x)+zg(x)-P(x)
   =(f(x)-P_f(x))+z(g(x)-P_g(x)).
\]
Both coefficients in the final expression lie in \(K\).  Since
\(z\notin K\), the elements \(1,z\) are linearly independent over \(K\).  Hence
both coefficients vanish for every \(x\in S\), proving the separate explanations
of \(f\) and \(g\).  Thus any explained slope outside \(K\) is support-wise
contained, so a noncontained slope must lie in \(K\).
\end{proof}

\begin{remark}[extension-field warning]
Theorem~\ref{thm:subfield-confinement-base-lines} is intentionally one-sided.
It only treats \(K\)-valued received lines.  It does not rule out genuinely
\(F\)-valued residue-line witnesses, and therefore must not be used as an
extension-field MCA theorem without an additional hypothesis forcing the line to
be \(K\)-valued.
\end{remark}

\begin{theorem}[seven-slope arithmetic gate over \(\F_{17^{32}}\)]
\label{thm:seven-slope-arithmetic-gate}
Let \(\qline=17^{32}\).  If a row over \(\F_{17^{32}}\) has
\(\operatorname{LD}_{\rm sw}(C,264)\ge7\), then at radius
\[
  \delta=1-\frac{264}{512}=\frac{31}{64}
\]
its support-wise MCA error is strictly larger than \(2^{-128}\).  In particular,
for the \(n=512,k=256\) row, any strict264 seven-slope certificate proves
\[
  \varepsilon_{\rm mca}(C,31/64)>2^{-128}.
\]
\end{theorem}

\begin{proof}
The MCA sampler over finite slopes has denominator \(\qline=17^{32}\).  Seven
bad slopes therefore give error at least \(7/17^{32}\).  The exact integer
comparison is
\[
  7\cdot 2^{128}-17^{32}
  =14064973686101998399515954701426232831>0.
\]
Thus \(7/17^{32}>2^{-128}\).  Agreement \(264\) on a domain of size \(512\)
corresponds to relative radius \(1-264/512=31/64\), giving the claim.
\end{proof}

\begin{theorem}[one-step staircase pinning criterion]
\label{thm:one-step-staircase}
Fix a finite-slope sampler of denominator \(\qline\) and a target
\(\varepsilon_*=2^{-128}\).  Let
\[
  B(a)=\operatorname{LD}_{\rm sw}(C,a)
\]
for the allowed line family, and let \(\delta_a=1-a/n\).  If, for some integer
\(a\ge k\),
\[
  B(a)>\varepsilon_*\qline
  \qquad\text{and}\qquad
  B(a+1)\le \varepsilon_*\qline,
\]
then the closed-ball agreement staircase is unsafe at radius \(\delta_a\) and
safe on the next open interval below it, including the grid radius
\(\delta_{a+1}\).  Equivalently, the largest grid radius certified safe is
\(\delta_{a+1}\), while the continuous supremum of certified safe radii is
\(\delta_a\) and is not attained under the closed endpoint convention.

Before the tangent staircase was isolated, the special case
\[
  B(264)\ge7,
  \qquad
  B(265)\le6
\]
would have pinned the finite row to an unsafe/safe transition between
\[
  \delta_{265}=247/512
  \qquad\text{and}\qquad
  \delta_{264}=31/64.
\]
For the \(\F_{17^{32}}\), \(n=512\), \(k=256\) row, the premise
\(B(265)\le6\) is false by the moving-root tangent floor.  The actual
finite-slope staircase gate supplied by Corollary~\ref{cor:f17-32-exact-tangent-gate}
is \(B(506)=7\), \(B(507)=6\).
\end{theorem}

\begin{proof}
The function \(B(a)\) is nonincreasing in \(a\), because raising the agreement
threshold only removes valid supports.  At radius \(\delta_a\), the closed-ball
predicate allows agreement at least \(a\), so the error is \(B(a)/\qline\), which
is larger than \(\varepsilon_*\) by assumption.  At radius \(\delta_{a+1}\), the
predicate requires agreement at least \(a+1\), so the error is
\(B(a+1)/\qline\), which is at most \(\varepsilon_*\).

More generally, any radius \(\delta<\delta_a\) sufficiently close to
\(\delta_a\) has integer distance bound at most \(n-a-1\), hence agreement at
least \(a+1\).  Thus the closed endpoint at \(\delta_a\) is the first unsafe
grid step, while the safe side is immediately below that endpoint.  For
\(\qline=17^{32}\), the inequalities
\[
  6\cdot2^{128}<17^{32}<7\cdot2^{128}
\]
convert the printed integer conditions into the two displayed error inequalities,
and \(1-265/512=247/512\), \(1-264/512=31/64\).
\end{proof}

\section{June 29 Finite-Slope and Interleaved-List PR Packets}
\label{sec:jun29-pr-packets}

This section records the mathematically useful parts of PRs \#131--\#135.  The
branches were not merged wholesale, because they were stale relative to the
current Paper D v6 and public-site state.  The integrated source files are the
compact notes and verifiers listed in
\path{experimental/notes/triage/pr-triage-2026-06-29.md}.

\subsection{Boundary-off external-anchor normal form}

PR \#131, due to AllenGrahamHart, contributes a local M1 proof-program normal
form.  For a shadow \(S\subset D\) of size \(j-1\) and an external anchor
\(\beta\notin D\), the one-outside locator satisfies
\[
  \ell_{S,\beta}=\ell_S^+-\beta\ell_S^0 .
\]
Therefore the Hankel landing matrix
\[
  M_S(\beta)=
  \bigl[H(u)\ell_{S,\beta}\quad H(v)\ell_{S,\beta}\bigr]
\]
has columns affine-linear in \(\beta\), and the rank-one condition is cut out
by quadratic minors in \(\beta\).  Thus, for fixed \(S\), either a nonzero
quadratic minor gives at most two external anchors, or all minors vanish
identically and the branch is a ruled rank-one pencil.  The later
same-slope/root-slice note now classifies this ruled \(t=2\) Hankel branch:
it is inactive or fixed finite slope, and the fixed-slope case lifts to a
boundary root-slice core \(H_{3,j-1}(u+z_0v)\ell_S=0\).  Therefore after
fixed-slope boundary slices are charged, each shadow has at most two active
non-ruled external anchors.  This is recorded as PROVED-LOCAL /
RULED-BRANCH CLASSIFIED / PROOF-PROGRAM / AUDIT, not as a proof of M1.

\subsection{\(a=507\) adjacent bridge route cut}

PR \#132, due to Scott Hughes, records a useful non-promotion result for the
\(\F_{17^{32}},n=512,k=256\) row.  The tangent-star theorem gives
\[
  \operatorname{LD}_{\rm sw}(C,a)=513-a
  \qquad (a\ge427),
\]
so in particular
\[
  \operatorname{LD}_{\rm sw}(C,507)=6.
\]
Any bridge that converted an adjacent line-plus-interleaved-list
``\(6+1=7\)'' ledger into one additional same-denominator finite-slope
support-wise event would imply
\(\operatorname{LD}_{\rm sw}(C,507)\ge7\), contradicting the exact
tangent-star bound.  Hence the adjacent arithmetic remains a separate
BRIDGE-NEEDED ledger and is not a new MCA row.

\subsection{Interleaved-list lower row at agreement \(326\)}

PR \#133, due to Scott Hughes, belongs to the interleaved-list track, not the
MCA slope track.  It proves high-agreement uniqueness
\[
  \Lambda_\mu(C,a)=1\qquad (384\le a\le512)
\]
and a support-occupancy route cut excluding seven witnesses for
\[
  373\le a\le383.
\]
It also gives lower-bound constructions at agreements \(291\), \(292\), \(320\),
and, most usefully for the board, a hybrid quotient-residual construction with
\[
  \Lambda_\mu(C,326)\ge7 .
\]
Since
\[
  6\cdot2^{128}<17^{32}<7\cdot2^{128},
\]
this clears the separate interleaved-list numerical gate over denominator
\(|F|=17^{32}\).  The packet does not claim an MCA numerator, protocol
soundness, ordinary list-decoding exactness, or an exact \(\delta_C^*\).

\subsection{Random simple-pole finite-slope floors}

PR \#134, due to Vadim Avdeev, proves a general random simple-pole lower bound.
For \(C=\RS[\F,D,k]\), \(|F|=q\), \(|D|=n\), \(\beta\in F\setminus D\), and
\(k<a\le n\), put
\[
  C_a=\binom{n}{a},\qquad s=a-k,\qquad
  R_a=\sum_{v=0}^{a-k-1}
        \binom{a}{v}\binom{n-a}{v}q^{-v}.
\]
Then there is a simple-pole line with at least
\[
  \left\lceil
    \frac{qC_a}{C_a+R_aq^s}
  \right\rceil
\]
finite bad slopes at agreement \(a\).  In the row
\[
  F=\F_{17}[z]/(z^{32}-3),\quad H=\langle z\rangle,\quad |H|=512,\quad k=256,
\]
the verified floors are
\[
\begin{array}{c|c}
a & \text{finite bad-slope lower bound}\\
\hline
257 & 17^{32}\\
258 & 17^{32}\\
259 & 17^{32}-68{,}904\\
260 & 33{,}439{,}260{,}151{,}101{,}646{,}297{,}506{,}087{,}371{,}119{,}470 .
\end{array}
\]
The direction is a genuine simple-pole direction: it is not a degree-\(<256\)
codeword on any support of size \(>256\).

\subsection{Coset-packet finite-slope floors for \(a=261,\ldots,288\)}

PR \#135, also due to Vadim Avdeev, continues the finite-slope floor beyond the
range where the random simple-pole second-moment bound is useful.  The
construction uses full dyadic cosets \(K_M\le H\), whose locators have the form
\[
  \prod_{x\in hK_M}(X-x)=X^M-h^M,
\]
so unions of full \(K_M\)-cosets share low elementary coefficients.  After
fixing a small partial coset when needed, this produces common-prefix support
families.  A separating deep anchor then turns the corresponding numerator
polynomials into distinct finite bad slopes.

For the same \(\F_{17^{32}},n=512,k=256\) row, the certified floors are
\[
\begin{array}{c|c}
a & \text{finite bad-slope lower bound}\\
\hline
261\ldots263 & \binom{63}{32}\\
264 & \binom{64}{33}\\
265\ldots271 & \binom{31}{16}\\
272 & \binom{32}{17}\\
273\ldots287 & \binom{15}{8}\\
288 & \binom{16}{9}.
\end{array}
\]
Every displayed row beats the moving-root tangent floor \(513-a\), but these
are still lower-bound rows only.  They do not prove a safe-side theorem, an
ordinary list-decoding theorem, an interleaved-list statement, or protocol
soundness.

\section{Non-Claims}
\label{sec:nonclaims}

This note does not claim:
\begin{itemize}
\item a proof of the corrected MCA conjecture;
\item a positive theorem above corrected reserve;
\item a \(\qgen\) local-limit theorem;
\item a list, CA, MCA, line-decoding, SNARK, or protocol consequence;
\item a large \(\Theta(p^2)\) counterpacket;
\item a new prize-worthy numerator beyond the current Cycle116/119
  \(52{,}747{,}567{,}092\) finite/source-scoped record;
\item priority for the BCHKS quotient-locator proof pattern.
\end{itemize}

The mathematical value is narrower.  On the imported side, the file gives a
provenance-safe dictionary and wrapper for using BCHKS without claiming it as
new.  On the local Paper B side, the possible large branch in the restricted
quadratic-extension window must satisfy an explicit algebraic identity before
it can even be surface-sized.

\appendix

\section{BibTeX Entry to Copy into Main Papers}

\begin{verbatim}
@misc{BCHKS25,
  author = {Eli Ben-Sasson and Dan Carmon and Ulrich Habock and
            Swastik Kopparty and Shubhangi Saraf},
  title  = {On Proximity Gaps for Reed--Solomon Codes},
  year   = {2025},
  note   = {Manuscript dated November 11, 2025},
  url    = {https://www.math.toronto.edu/swastik/rs-proximity-gaps-2025.pdf}
}
\end{verbatim}

\begin{thebibliography}{99}
\bibitem{ProximityPrize26}
The Proximity Prize.  \emph{The Proximity Prize: Reed--Solomon Challenge}.
Preliminary version, accessed June 26, 2026.
\url{https://proximityprize.org/}.

\bibitem{BCHKS25}
Eli Ben-Sasson, Dan Carmon, Ulrich Hab\"ock, Swastik Kopparty, and Shubhangi
Saraf.  \emph{On Proximity Gaps for Reed--Solomon Codes}.  Manuscript dated
November 11, 2025.
\url{https://www.math.toronto.edu/swastik/rs-proximity-gaps-2025.pdf}.
\end{thebibliography}

\end{document}
